% ConsistenCy v3 — Theoretical Foundations, Mathematical Formalization & Research Report
% XeLaTeX + biblatex/biber.  Build: latexmk -xelatex -outdir=build main.tex
\documentclass[UTF8,zihao=-4,a4paper,oneside,openany,fontset=windows]{ctexbook}

\usepackage[a4paper,left=2.6cm,right=2.6cm,top=2.8cm,bottom=2.8cm]{geometry}
\usepackage{amsmath,amssymb,amsthm,mathtools,bm}
\usepackage{tikz}
\usetikzlibrary{arrows.meta,positioning,shapes.geometric,shapes.misc,calc,fit,backgrounds,decorations.pathreplacing}
\usepackage{pgfplots}
\pgfplotsset{compat=1.18}
\usepackage{booktabs,longtable,tabularx,array,makecell}
\usepackage{enumitem}
\usepackage{caption}
\usepackage[table]{xcolor}
\usepackage{fancyhdr}
\usepackage{adjustbox}
\usepackage{url}
\setlength{\headheight}{13pt}
\setlength{\emergencystretch}{3em}
\usepackage[backend=biber,style=numeric-comp,sorting=none,maxbibnames=9,minbibnames=3,doi=true,isbn=false,url=false,eprint=true]{biblatex}
\addbibresource{references.bib}

\definecolor{linkblue}{RGB}{20,70,150}
\definecolor{citegreen}{RGB}{0,110,110}
\definecolor{inkgray}{RGB}{80,80,80}
\definecolor{boxgray}{RGB}{242,244,247}
\definecolor{framegray}{RGB}{160,168,178}
\definecolor{accentred}{RGB}{170,50,45}
\definecolor{accentblue}{RGB}{50,90,170}
\definecolor{softgreen}{RGB}{235,246,238}
\definecolor{softblue}{RGB}{235,242,250}
\definecolor{softamber}{RGB}{250,244,230}
\definecolor{softrose}{RGB}{250,236,236}

\usepackage{hyperref}
\hypersetup{colorlinks=true,linkcolor=linkblue,citecolor=citegreen,urlcolor=linkblue,
  pdftitle={ConsistenCy v3 理论基础、数学形式化与研究报告},
  pdfauthor={ConsistenCy Research},
  pdfsubject={Theory \& Foundations Technical Report},
  pdfcreator={XeLaTeX + biblatex}}

% ---------- theorem environments ----------
\theoremstyle{plain}
\newtheorem{theorem}{定理}[chapter]
\newtheorem{proposition}[theorem]{命题}
\newtheorem{corollary}[theorem]{推论}
\theoremstyle{definition}
\newtheorem{definition}[theorem]{定义}
\newtheorem{openproblem}[theorem]{开放问题}
\theoremstyle{remark}
\newtheorem{remark}{注}[chapter]
% assumptions numbered globally as printed labels A1, A2, ...
\newcounter{assump}
\renewcommand{\theassump}{A\arabic{assump}}
\theoremstyle{definition}
\newtheorem{assumptioninner}[assump]{假设}
% classification tag shown at the right of proposition statements
\newcommand{\tagm}[1]{\hfill{\small\itshape\color{inkgray}[#1]}}
\newcommand{\DIRECT}{\textsc{Direct}}
\newcommand{\STRONG}{\textsc{Strong}}
\newcommand{\USEFUL}{\textsc{Useful}}
\newcommand{\PROPOSED}{\textsc{Proposed}}
\newcommand{\WEAK}{\textsc{Weak}}
\newcommand{\NA}{\textsc{N/A}}

% ---------- math macros ----------
\newcommand{\Auth}{\operatorname{Auth}}
\newcommand{\Exec}{\operatorname{Exec}}
\newcommand{\Avail}{\operatorname{Available}}
\newcommand{\Vis}{\operatorname{Visible}}
\newcommand{\ground}{\operatorname{ground}}
\newcommand{\fp}{\operatorname{fp}}
\newcommand{\canon}{\operatorname{canonical}}
\DeclareMathOperator{\E}{\mathbb{E}}
\DeclareMathOperator{\Var}{\mathbb{V}\mathrm{ar}}
\newcommand{\Nat}{\mathbb{N}}
\newcommand{\States}{S_{\mathrm{ACB}}}
\newcommand{\tcode}[1]{\texttt{#1}}

\captionsetup{font=small,labelfont=bf}
\setlength{\parskip}{0.35em}

% ---------- header / footer ----------
\pagestyle{fancy}
\fancyhf{}
\fancyhead[L]{\small\itshape ConsistenCy v3 理论基础研究报告}
\fancyhead[R]{\small\nouppercase\leftmark}
\fancyfoot[C]{\thepage}
\renewcommand{\headrulewidth}{0.4pt}

\begin{document}

% ================= cover =================
\begin{titlepage}
\centering
\vspace*{2.2cm}
{\color{inkgray}\small THEORY \& FOUNDATIONS TECHNICAL REPORT\par}
\vspace{0.8cm}
{\color{framegray}\rule{\textwidth}{1.2pt}\par}
\vspace{1.2cm}
{\Huge\bfseries ConsistenCy v3\\[0.5em] 理论基础、数学形式化与研究报告\par}
\vspace{1.0cm}
{\Large Theoretical Foundations, Mathematical Formalization\par}
{\Large and a Research Agenda\par}
\vspace{1.4cm}
{\color{framegray}\rule{0.618\textwidth}{0.6pt}\par}
\vspace{1.2cm}
{\large 基于 Capability Security、Queueing Theory、Provenance、\par}
{\large Agent Systems 与 Information Freshness 文献的诚实映射\par}
\vspace{2.4cm}
\begin{minipage}{0.8\textwidth}\centering\small
\begin{tabular}{rl}
报告类型 & 理论与技术报告（非产品文档）\\
研究对象 & ConsistenCy v3（分支 \tcode{v3}，工作树状态 2026-08-28）\\
文献基线 & 45 篇经网络核验的文献（20 篇承重核心）\\
版本 & v1.0\\
日期 & 2026 年 8 月 28 日\\
\end{tabular}
\end{minipage}
\vfill
{\color{inkgray}\small 本报告不修改任何产品源代码；所有架构事实以当前源码为唯一真相来源。\par}
\end{titlepage}

\frontmatter
% ---------- abstract ----------
\begin{center}
{\Large\bfseries 摘要}
\end{center}
\noindent
本报告对 ConsistenCy v3——一个仓库原生的 Agent Harness 操作系统——进行理论基础研究。以当前工作树（分支 \tcode{v3}，2026-08-28）为唯一事实来源，我们把 14 个经源码核验的组件映射到能力安全、排队与调度、溯源与可复现性、Agent 系统与工作流形式方法、信息新鲜度五个理论域，并构造共享记号层面的形式模型 $\mathcal{M}_C$：全局状态 $X_t$（十元组，见附录~\ref{app:symbols}）、九步授权谓词 $\Auth$、撤销支配命题、可用性/可见性与权威的正交性命题、ACB 与 Fiber 状态机的单向耦合关系、工作流 DAG 的裁决面独立性、$M/M/c$ 基线排队模型（显式声明的近似）、证据指纹的确定性与碰撞界、以及版本滞后度量 $D(t)$ 与检测时延界。所有映射按六类支持强度分级（DIRECT、STRONG、USEFUL、PROPOSED、WEAK、NA），全部数学结论按五类标签分类并附假设、证明与失败条件；三组负结果（Petri 网机件、CUSUM 变化检测、线性时间年龄）被显式记录。文献基线为 45 篇经网络核验的文献（20 篇承重核心）；模型经过敌意数学审计（全部 FATAL/MAJOR 修正已并入）与独立引文审计（零虚构）。核心结论：ConsistenCy 的安全内核在\textbf{机制层}是引用监视器与对象能力理论的直系实例化（\DIRECT；分析性与验证性差距见第 19 章），其“Coeffect $\neq$ Authorization”“Context VM $\neq$ Permission System”公理获得精确数学表述——安全只由 $\Auth$ 承担；而调度优化、新鲜度度量与模型检验被诚实标记为提议研究方向，不构成现状声明。

\medskip
\noindent\textbf{关键词：}能力安全；引用监视器；Agent Harness；形式化验证；排队论；数据溯源；信息新鲜度；LLM Agent 治理
\clearpage

\tableofcontents
\listoffigures
\listoftables
\clearpage

\mainmatter

% 第 1-2 章：项目定位与研究动机
\chapter{ConsistenCy v3 是什么}\label{ch:intro}

\section{一句话定位}

ConsistenCy v3 是一个\textbf{仓库原生（Repository-Native）的 Agent Harness 操作系统}：它以内核化的方式为“多 Agent 持续修改同一仓库”这一场景提供权限、上下文、调度、证据与运行生命周期的统一治理，而\emph{公共 PR 审查（Public PR Review）只是其上运行的第一个工作负载}。系统的物理核心可以写成一个极简公式：
\begin{equation}
\text{ConsistenCy v3} \;=\; \text{Kernel} \;+\; \text{Cordis Harness} \;+\; \text{Evidence Engine},
\end{equation}
其中 Kernel 负责能力授权（Capability Broker）、系统调用中介（Syscall Gateway）、调度（Kernel Scheduler）、上下文虚拟机（Context VM）与审计日志（Audit Journal）；Cordis Harness 提供响应式的 Fiber 生命周期与共效（Coeffect）依赖注入；Evidence Engine 以确定性分析器与内容寻址指纹提供可复现的证据底座。

本报告的研究对象是\textbf{当前工作树中的真实实现}（分支 \tcode{v3}，2026-08-28，含未提交的 CKPT5 变更），而非任何历史设计稿。所有架构事实均以只读源码调查为准，并给出文件与行号定位；文档与代码不一致之处（例如能力授权检查“8 步”注释与实际 9 步实现的差异）一律以代码为准并在正文中如实记录。

\section{产品要回答的五个问题}

ConsistenCy v3 的产品核心价值可以概括为五个问题：
\begin{enumerate}[itemsep=0.1em]
  \item \textbf{谁被允许做什么}（Who may do what?）——Kernel Capability 授权；
  \item \textbf{哪些组件当前应当存活}（Which components should be alive?）——Cordis Fiber 与共效生命周期；
  \item \textbf{谁现在运行}（Who runs now?）——Kernel Scheduler 的准入、并发与预算；
  \item \textbf{Agent 看见什么}（What does the agent see?）——Context VM 的 ContextImage 与工作集；
  \item \textbf{什么事实证明结果}（What facts prove the result?）——Evidence Engine、RepositorySnapshot 与溯源。
\end{enumerate}
这五个问题恰好构成本报告五个理论主题（能力安全、生命周期与授权分离、可见性与权限分离、调度与资源治理、证据与可复现性）的工程原型。换言之，本报告的任务不是给项目“贴数学”，而是把工程上已经成立的边界条件翻译成可以被检查、讨论与继续研究的数学对象。

\section{本报告的证据纪律}

本报告执行严格的分类纪律，任何“理论映射”都必须归入下列六类之一，并在正文与附录 B 的映射矩阵中显式标注：
\begin{itemize}[itemsep=0.05em]
  \item \DIRECT{}（直接理论支撑）：该机制是已有理论的标准实例化；
  \item \STRONG{}（强类比）：理论结构与机制同构，但场景假设不完全重合；
  \item \USEFUL{}（有用的形式化）：理论提供描述工具而非结论；
  \item \PROPOSED{}（提议的形式化/扩展）：本报告首次给出的模型，非已确立理论；
  \item \WEAK{}（弱或投机）：仅作背景，不承担论证责任；
  \item \NA{}（不适用）：经研究判定不应映射，作为负结果记录。
\end{itemize}
所有数学结论同样分类为：\textsc{TheoremFromKnownModel}（已知模型的定理实例）、\textsc{Proposition}（ConsistenCy 模型级命题）、\textsc{EngineeringInvariant}（工程不变式，源码可证）、\textsc{Conjecture}（猜想）与\textsc{ProposedMetric}（提议度量）。除非明确声明，任何命题证明的都是\textbf{架构模型} $\mathcal{M}_C$ 的性质，而不是运行中实现的性质；源码可证的部分会单独标注。报告还区分三类贡献：\textbf{既有工程贡献}（项目已实现）、\textbf{形式化解释}（本报告首次给出）与\textbf{新研究假设}（未来可验证）。

\chapter{为什么 Agent Harness 需要理论基础}\label{ch:why}

\section{问题背景：LLM Agent 的结构性风险}

大模型 Agent 系统在 2023--2026 年间快速发展，但主流框架（如以 ReAct~\cite{yao2023react} 为代表的推理-行动循环、以 AutoGen~\cite{wu2024autogen} 为代表的多 Agent 会话编程、以 MetaGPT~\cite{hong2024metagpt} 为代表的 SOP 编排工作流）普遍将\textbf{权限}问题留给宿主应用的约定，而非系统的构造性属性。经验研究也印证了这一风险的现实性：对大规模真实代码审查评论的实证测量表明，LLM 生成的审查评论中被开发者实际处置的比例呈现强烈的类型依赖性，原始模型输出本身并不是可靠的审查工件~\cite{goldman2025}。攻击面研究进一步表明，提示注入（prompt injection）本质上是\emph{混淆代理人}（confused deputy）问题在 LLM 时代的再现~\cite{hardy1988,debenedetti2025}：拥有权限的代理人被不可信内容诱导，以其\emph{环境固有权限}（ambient authority）执行违背委托人意图的操作。

ConsistenCy 的核心论点是：\textbf{当 Agent 成为仓库的一等修改者时，Agent Harness 必须像一个操作系统内核一样回答权限问题}，而这类问题恰恰是系统安全领域六十余年积累的成熟理论对象——能力系统（capability-based security）~\cite{dennis1966,saltzer1975,miller2003,miller2006}、引用监视器（reference monitor）~\cite{anderson1972}、撤销理论~\cite{redell1974}。

\section{工程直觉为何不够}

ConsistenCy v3 的七条架构公理（Coeffect $\neq$ Authorization；Effect $\neq$ Rollback；Ring $\neq$ Execution Domain；Cordis Context Isolation $\neq$ Process Security Isolation；Context VM $\neq$ Permission System；Scheduler 属于 Kernel；Evidence 是 correctness 而非日志）目前以工程断言的形式存在。它们在实践中有效，但存在三个未被回答的问题：
\begin{enumerate}[itemsep=0.1em]
  \item \textbf{可检查性}：公理之间是否一致？“撤销在下一个系统调用处同步生效”与“Fiber 生命周期异步传播”并存的正确性条件是什么？
  \item \textbf{可比较性}：缺少形式化描述时，无法将 ConsistenCy 与 seL4~\cite{klein2009}、Capsicum~\cite{watson2010}、in-toto~\cite{torresarias2019} 等系统在\emph{同一维度}上比较，只能停留在术语层面的类比。
  \item \textbf{可研究性}：调度、新鲜度、证据接地等机制的未来改进（例如预算感知准入、基于版本滞后的触发策略）需要数学模型才能提出并证伪。
\end{enumerate}

\section{本报告的立场}

本报告的立场可以概括为一句话：
\begin{quote}
\itshape 找到 ConsistenCy 的真实理论基础，并把项目已有的工程直觉转化为可以被检查、讨论和继续研究的数学对象；同时诚实标记所有“看似相关实则不然”的理论，把它们作为负结果记录。
\end{quote}
因此，报告包含大量\textbf{非主张}（non-claims）：当前调度器\emph{不是}约束马尔可夫决策过程（CMDP）优化器；信息年龄（Age of Information）对仓库审查新鲜度\emph{不是}正确的原始度量；当前架构\emph{没有}经过形式化验证（seL4 意义上）；证据接地在遗留评审 schema 层允许无证据引用的 finding 存在。这些非主张与正面结果同等重要——它们划定了“已确立理论”与“提议模型”的边界，防止理论包装（theory washing）。

% 第 3 章：当前系统架构与形式化对象
\chapter{当前系统架构与形式化对象}\label{ch:arch}

本章把当前工作树中的 ConsistenCy v3 架构压缩为后续形式化所需的组件清单。每个组件的“工程不变式”都经过只读源码核验（定位到文件与行号），第~\ref{ch:security} 章起再逐一给出数学对象。

\section{三层物理结构与形式化对象总览}

图~\ref{fig:theorymap} 给出“架构组件 $\rightarrow$ 理论域”的总映射。每个组件旁标注了本报告采用的主要理论域与支持强度（六类分级见第~\ref{ch:intro} 章）；详细的逐组件论证分布于第~\ref{ch:security}--\ref{ch:freshness} 章，完整的文献映射矩阵见附录~\ref{app:mapping}。

\begin{figure}[htbp]
\centering
\begin{adjustbox}{max width=\textwidth}
\begin{tikzpicture}[
  comp/.style={draw=framegray,fill=boxgray,rounded corners=2pt,minimum height=1.55em,inner sep=4pt,font=\scriptsize,align=center},
  dom/.style={draw=accentblue!60,fill=softblue,rounded corners=2pt,minimum height=1.55em,inner sep=4pt,font=\scriptsize\itshape,align=center},
  lnk/.style={-{Stealth[length=2mm]},draw=inkgray!70},
  lab/.style={font=\tiny\itshape\color{inkgray},align=center}]
% Kernel column
\node[comp] (broker) at (0,0) {CapabilityBroker\\（9 步授权链）};
\node[comp] (gateway) at (0,-1.25) {SyscallGateway\\+ CommitCoordinator};
\node[comp] (sched) at (0,-2.5) {KernelScheduler\\（3 路准入）};
\node[comp] (context) at (0,-3.75) {Context VM\\（页面表/COW）};
\node[comp] (audit) at (0,-5.0) {AuditJournal\\（仅追加）};
\node[comp,fit=(broker)(gateway)(sched)(context)(audit),inner sep=7pt,fill=softblue!35,draw=accentblue!40] (kernelbox) {};
\node[font=\small\bfseries\color{accentblue}] at (kernelbox.north) [above=1pt] {Kernel 层};
% Harness
\node[comp] (fiber) at (4.4,-1.25) {Cordis Fiber\\生命周期};
\node[comp] (coe) at (4.4,-2.5) {Coeffect\\声明式依赖};
\node[comp,fit=(fiber)(coe),inner sep=7pt,fill=softgreen!50,draw=framegray] (harnessbox) {};
\node[font=\small\bfseries] at (harnessbox.north) [above=1pt] {Cordis 层};
% Workload
\node[comp] (sup) at (4.4,-4.3) {Supervisor Planner\\+ 评审 Agent};
\node[comp] (wf) at (4.4,-5.5) {Workflow Runtime\\（DAG/触发）};
% Evidence
\node[comp] (evid) at (8.8,-1.25) {EvidenceStore\\（内容指纹）};
\node[comp] (snap) at (8.8,-2.5) {RepositorySnapshot\\（SHA 钉定）};
\node[comp] (eng) at (8.8,-3.75) {确定性分析器\\（Python/TS）};
\node[comp,fit=(evid)(snap)(eng),inner sep=7pt,fill=softamber!60,draw=framegray] (evbox) {};
\node[font=\small\bfseries] at (evbox.north) [above=1pt] {Evidence 层};
% Theory domains (right)
\node[dom] (t1) at (13.0,0) {能力安全 / 引用监视器\\（\DIRECT）};
\node[dom] (t2) at (13.0,-1.25) {排队的业务效果建模\\（幂等意图）};
\node[dom] (t3) at (13.0,-2.5) {排队论 / 准入控制\\（\DIRECT{} 词汇）};
\node[dom] (t4) at (13.0,-3.75) {信息可见性 $\neq$ 权限\\（\DIRECT{} 分离）};
\node[dom] (t5) at (13.0,-5.0) {防篡改日志 / 事件溯源\\（\STRONG）};
\node[dom] (t6) at (13.0,-6.3) {监督式生命周期\\（\STRONG）};
\node[dom] (t7) at (13.0,-7.6) {工作流 DAG / 数据流\\（\STRONG）};
\node[dom] (t8) at (13.0,-8.9) {溯源 / 可复现计算\\（\STRONG--\DIRECT）};
\draw[lnk] (broker.east) to[bend left=8] (t1.west);
\draw[lnk] (gateway.east) to[bend left=4] (t2.west);
\draw[lnk] (sched.east) to[bend left=2] (t3.west);
\draw[lnk] (context.east) to (t4.west);
\draw[lnk] (audit.east) to[bend right=2] (t5.west);
\draw[lnk] (harnessbox.east) to[bend left=6] (t6.west);
\draw[lnk] (wf.east) to (t7.west);
\draw[lnk] (evbox.east) to[bend right=6] (t8.west);
\end{tikzpicture}
\end{adjustbox}
\caption{架构 $\rightarrow$ 理论域总映射（ConsistenCy Architecture--Theory Domain Map）。支持强度为六类分级的简写；每条映射的完整论证与关键文献见对应章节及附录 B。}
\label{fig:theorymap}
\end{figure}

\section{组件清单与工程不变式}

下表汇总源码核验后的 14 个组件及其当前不变式（详细档案见研究笔记；行号针对 2026-08-28 工作树）：

\begin{longtable}{@{}>{\raggedright\arraybackslash}p{0.21\textwidth}>{\raggedright\arraybackslash}p{0.44\textwidth}>{\raggedright\arraybackslash}p{0.27\textwidth}@{}}
\toprule
\textbf{组件} & \textbf{源码核验的不变式（摘要）} & \textbf{主要定位}\\
\midrule
\endhead
CapabilityBroker & 授权 = 9 个谓词的有序首败拒绝链，每次受保护系统调用都重新求值；撤销/过期单调；允许与拒绝都写审计 & \path{capability/broker.ts:245-328}\\
SyscallGateway & \tcode{intent} 类动作（\tcode{repo.write}、\tcode{github.publish}）永不内联分发，必须经 CommitCoordinator；\tcode{llm.invoke} 为 commit/direct，用量由可信 Ring-1 处理器提交 & \path{syscall/authorize.ts:46-102}\\
CommitCoordinator & 幂等键去重的意图接受器；$payloadHash=\mathrm{SHA256}(\canon(payload))$；内存态，持久性委托给注入的 Sink & \path{commit/coordinator.ts:39-183}\\
KernelScheduler & 3 路准入（截止期取消 $\rightarrow$ 容量 $\rightarrow$ 优先级-先来先服务，严格平局决断 = 确定性）；终态吸收 & \path{scheduler/scheduler.ts:225-271}\\
AgentControlBlock & 12 状态 LTS；能力引用只存 12 位十六进制\emph{指纹}（描述符，非凭证） & \path{agent/state.ts:27-50}\\
Cordis Fiber & PENDING / LOADING / ACTIVE / UNLOADING（Cordis 4.0.0-rc.8 库语义）；共效缺失 $\Rightarrow$ PENDING；撤销经微任务异步传播 & \path{harness-core/runtime/adapter.ts:46-130}\\
Context VM & 页面不可变、SHA-256 内容寻址；驻留级在覆盖层；\tcode{pinned$\rightarrow$evicted} 抛错（失败关闭）；页面持有不授予任何权限 & \path{context/manager.ts:16-19,102-174}\\
AuditJournal & 仅 \tcode{record}+\tcode{entries} 接口（无可变性 API）；5 类事件；只出现指纹与哈希，永不出现原始句柄 & \path{audit/types.ts:21-100}\\
Evidence Engine & 指纹 $=$ SHA-256（11 元组规范 JSON）；\tcode{provenance.sha} 参与哈希；store 自行计算指纹、不信任分析器 & \path{evidence/fingerprint.ts:35-103}\\
RepositorySnapshot & 以 \tcode{rev-parse --verify} 失败关闭地钉定 \tcode{headSha}；读取只经对象库；后续工作树变更不影响既有快照 & \path{snapshot/snapshot.ts:1-120}\\
Workflow Runtime & 校验（Kahn）$\rightarrow$ 编译（注册表解析，不发放持久授权）$\rightarrow$ 运行（逐系统调用授权）；触发是“纯观测溯源，绝非授权” & \path{workflow-runtime/validate,compile,host.ts}\\
触发计划 & 持久去重键 $\mathrm{SHA256}[\text{domain},\text{repo},\text{def},\text{event}]$；围栏式单飞执行；单次尝试，无重试循环 & \path{workflow-runtime/triggers.ts,store.ts}\\
仓库监督 & 轮询 + 观测摘要去重 + 防抖；每个 (repo, head, 配置摘要) 至多一个事件；只观测记录，不自行入队 & \path{heartbeat/repositorySupervisor.ts}\\
LLM 网关 & Ring-3 到提供商的唯一路径是逐调用授权的 \tcode{llm.invoke}；密钥仅服务端；预算是唯一在代码内的反压 & \path{workload-review/facades/llm-facade.ts:73-140}\\
\bottomrule
\caption{组件级工程不变式（源码核验摘要）。}
\label{tab:components}
\end{longtable}

\section{权威结构图}

图~\ref{fig:authority} 展示系统的\textbf{权威（authority）流向}：所有外部效果必须穿过 Kernel 的授权谓词；Cordis 与工作流层只能“触发与描述”；证据层只提供可复现事实。该图是后续所有安全命题（S1--S3）的几何底图。

\begin{figure}[htbp]
\centering
\begin{adjustbox}{max width=\textwidth}
\begin{tikzpicture}[
  layer/.style={draw=framegray,rounded corners=3pt,inner sep=6pt},
  box/.style={draw=framegray,fill=boxgray,rounded corners=2pt,font=\scriptsize,align=center,minimum height=1.5em,inner sep=4pt},
  gate/.style={draw=accentred!80,fill=softrose,rounded corners=2pt,font=\scriptsize\bfseries,align=center,inner sep=5pt},
  lnk/.style={-{Stealth[length=2.2mm]},draw=inkgray!80,thick},
  nolnk/.style={draw=inkgray!40,dashed}]
\node[layer,fill=softblue!30] (api) at (0,4.2) {\shortstack[c]{\small\bfseries 触发与描述层（无权威）\\ Web UI / API / Workflow 定义 / 绑定 / 触发事件}};
\node[layer,fill=softgreen!40] (harness) at (0,2.5) {\shortstack[c]{\small\bfseries Cordis Harness（生命周期权威）\\ Fiber PENDING/ACTIVE $\cdot$ 共效可用性 $\cdot$ 异步传播}};
\node[gate] (authz) at (0,0.9) {\shortstack[c]{Kernel 授权谓词 $\Auth(a,\alpha,r,s,t)$\\ 9 步首败拒绝链 $\cdot$ 每次调用重评 $\cdot$ 失败关闭}};
\node[layer,fill=softamber!40] (eff) at (0,-0.9) {\shortstack[c]{\small\bfseries 效果世界\\ git 读 / LLM 调用 / GitHub 发布（经意图门）}};
\node[layer,fill=boxgray] (ev) at (6.4,0.9) {\shortstack[c]{\small\bfseries Evidence Engine（事实权威）\\ 快照 SHA 钉定 $\cdot$ 内容指纹 $\cdot$ 确定性分析器}};
\node[layer,fill=boxgray] (audit2) at (6.4,-0.9) {\shortstack[c]{\small\bfseries AuditJournal（历史权威）\\ 仅追加 $\cdot$ 全决策记录 $\cdot$ 指纹化}};
\draw[lnk] (api) -- (harness) node[midway,right=1pt,font=\tiny\itshape\color{inkgray}] {触发};
\draw[lnk] (harness) -- (authz) node[midway,right=1pt,font=\tiny\itshape\color{inkgray}] {受保护 syscall};
\draw[lnk] (authz) -- (eff) node[midway,right=1pt,font=\tiny\itshape\color{inkgray}] {允许才执行};
\draw[lnk] (authz) -- (audit2) node[midway,above,font=\tiny\itshape\color{inkgray}] {全决策入账};
\draw[lnk] (eff) -- (ev) node[midway,right=1pt,font=\tiny\itshape\color{inkgray}] {证据落地};
\draw[nolnk] (api.west) to[bend left=25] node[midway,left=2pt,font=\tiny\itshape\color{inkgray},align=right] {无直达路径\\（S3：数据非权威）} (eff.west);
\end{tikzpicture}
\end{adjustbox}
\caption{Kernel / Cordis / Agent / Evidence 权威结构图。实线为允许的权威路径；虚线标注被架构排除的路径（触发与描述层数据不能直接成为权威来源，详见命题~\ref{prop:verdictplane}）。}
\label{fig:authority}
\end{figure}

\section{形式化对象预览}

基于上述事实，本报告在后续章节构造以下数学对象（完整定义见第~\ref{ch:security}--\ref{ch:unified} 章，符号表见附录~\ref{app:symbols}）：
\begin{itemize}[itemsep=0.05em]
  \item 全局状态 $X_t=(R_t,S_t,Q_t,\mathfrak{A}_t,C_t,V_t,K_t,E_t,P_t,B_t)$ 与演化 $X_{t+1}=\Phi(X_t,u_t,\zeta_t)$（\PROPOSED{} 统一模型）；
  \item 授权谓词 $\Auth$ 与安全不变式 S1--S3（时序逻辑表述）；
  \item 可用性 $\Avail$、可见性 $\Vis$ 与 $\Auth$ 的正交性命题；
  \item 12 状态 ACB LTS 与 4 状态 Fiber LTS 的单向耦合关系；
  \item 工作流 DAG $G=(V,E)$、拓扑序 $\sigma$ 与“校验 $\neq$ 授权”的裁决面独立性；
  \item $M/M/c$ 基线排队模型（显式声明的近似）与 CMDP 研究模型（\PROPOSED）；
  \item 证据指纹函数 $\fp$、接地分类器 $\ground$ 与确定性/碰撞界命题；
  \item 版本滞后（Version Staleness）度量 $D(t)$ 与检测时延界（\PROPOSED{} 度量 + \STRONG{} 结构类比）。
\end{itemize}

% 第 4 章：Capability Security 与 Kernel Authority
\chapter{Capability Security 与 Kernel Authority}\label{ch:security}

\section{理论谱系}

ConsistenCy 的能力内核站在一条清晰的理论谱系上。Dennis 与 Van Horn 在 1966 年定义了\textbf{能力表（capability list）}：每个保护域持有一张不可伪造的引用表，每次访问都经内核检查~\cite{dennis1966}——这正是 CapabilityBroker 中“能力记录表 + 每次系统调用求值”的直接祖先。Anderson 报告（1972）提出\textbf{引用监视器（reference monitor）}的三条性质：总是被调用（完全中介）、防篡改、小到可以分析与测试~\cite{anderson1972}；Saltzer 与 Schroeder（1975）随后系统化了最小特权、完全中介、失败默认与机制经济性等设计原则~\cite{saltzer1975}。Hardy 的\textbf{混淆代理人}（1988）刻画了环境固有权限的危险：被诱导的特权代理人以委托人权限行事~\cite{hardy1988}——LLM Agent 的提示注入威胁正是该问题的现代实例~\cite{debenedetti2025}。Miller 等人拆除了能力系统的三个经典迷思（能力等价于 ACL、不可约束、不可撤销）~\cite{miller2003}，其博士论文进一步把对象能力模型与并发控制统一为“健壮组合”理论，其中 vat（事件循环受限的消息传递并发）模型与 ConsistenCy 的“事件循环串行化内核 + 响应式 Fiber”结构高度同构~\cite{miller2006}。撤销理论方面，Redell 与 Fabry 证明了\textbf{经中介间接引用}可实现选择性撤销~\cite{redell1974}；ConsistenCy 的内核记录表正是这个“中介点”，只是它选择了更强的语义——永久单调撤销。在现代系统一侧，seL4 给出了带完整函数正确性证明的能力内核~\cite{klein2009}，Capsicum 展示了能力模式可以在既有通用环境中落地~\cite{watson2010}；二者分别是 ConsistenCy “形似而未证”的金标准与实践先例。

\paragraph{映射强度判定。}
\begin{itemize}[itemsep=0.05em]
  \item “每次受保护调用重新授权 + 默认拒绝 + 单一权威点”$\rightarrow$ 引用监视器/完全中介：\DIRECT{}（机制类的标准实例化）；
  \item “不可猜测的 256 位句柄 + 指纹化审计”$\rightarrow$ 对象能力的不可伪造引用：\DIRECT{}（以熵而非硬件标签实现不可伪造性，弱于 CHERI 式架构保证，此处如实降格）；
  \item “永久单调撤销 + 下一调用同步拒绝”$\rightarrow$ 撤销经中介：\DIRECT{}（Redell--Fabry 问题的退化情形——不需要选择性，因此更易）；
  \item “LLM 作为不可信提议者”$\rightarrow$ 混淆代理人/POLA：\DIRECT{}（动机层）；CaMeL 式“能力即数据流控制”为平行现代工作（\STRONG，未经评审）~\cite{debenedetti2025}；
  \item “两阶段预算预留（reserve/commit）”$\rightarrow$ \emph{无}已确立理论（诚实的理论空白；最接近的已发表平行工作是运行时特权策略层）。
\end{itemize}

\section{授权谓词}\label{sec:auth}

\begin{definition}[授权谓词]\label{def:auth}
固定主语 $a$（Agent 身份 $(a.id,a.runId)$）、动作 $\alpha$、资源 $r$（带类别 $\kappa(r)$）、请求作用域 $s=(sha?,paths?)$ 与能力句柄 $h$。系统调用请求为四元组 $q=(h,a,\alpha,r,s)$。定义九个检查：
\begin{align}
\chi_1&:\; h\in\operatorname{dom}(C_t) & &\text{（句柄存在）}\nonumber\\
\chi_2&:\; \neg\, C_t(h).revoked & &\text{（未撤销）}\nonumber\\
\chi_3&:\; C_t(h).expiresAt > t & &\text{（未过期）}\nonumber\\
\chi_4&:\; C_t(h).subject = a.id & &\text{（主语匹配）}\nonumber\\
\chi_5&:\; \kappa(r)\in\{\text{evidence},\text{workspace}\} \Rightarrow C_t(h).runId = a.runId & &\text{（运行匹配）}\nonumber\\
\chi_6&:\; C_t(h).action = \alpha & &\text{（动作匹配）}\nonumber\\
\chi_7&:\; \operatorname{matchResource}(C_t(h).resource, r) & &\text{（按类别资源匹配）}\nonumber\\
\chi_8&:\; \operatorname{checkScope}(C_t(h).scope, s) & &\text{（作用域匹配）}\nonumber\\
\chi_9&:\; \operatorname{Reserve}_{B_t}(C_t(h),\Delta(\alpha))=\mathsf{ok} & &\text{（预算预留）}\nonumber
\end{align}
其中 $\chi_7$ 按类别定义（repository 按 id；snapshot 按 repositoryId+sha；evidence/workspace 按 runId；github.publish 可选钉定 pullNumber；llm 按 provider；ast 按 snapshotId）；$\chi_8$ 中“作用域 sha 已设而请求未带同等 sha”是\emph{违规}而非豁免，路径经 glob 匹配。授权谓词为有序首败拒绝链
\begin{equation}
\Auth(a,\alpha,r,s,t) \;\equiv\; \bigwedge_{i=1}^{9}\chi_i,
\end{equation}
在首个失败处返回 $(0,reason_i)$；允许与拒绝\textbf{都}写入审计日志。
\end{definition}

受保护动作集合 $\Pi$ 共 12 个注册项，按效果类 $\times$ 分发策略划分：pure/direct $\{\tcode{ast.query}\}$；read/direct 六项；revertible/direct $\{\tcode{workspace.write},\tcode{evidence.write}\}$；commit/direct $\{\tcode{llm.invoke}\}$；commit/intent $\Pi_{\mathrm{intent}}=\{\tcode{repo.write},\tcode{github.publish}\}$——后者永不内联分发，只能经 CommitCoordinator 的幂等意图门。

\begin{assumptioninner}[{完全中介（broker 层）}]\label{ass:a1}
$\Pi$ 中每个效果都经 \tcode{gateway.invoke} 到达 \tcode{broker.authorise}。\textbf{边界条件}：网关对\emph{未注册}动作是穿透的（网关层无显式未知系统调用拒绝）；中介完备性依赖 broker 的 $\chi_1$ 与发放策略拒绝——发放校验 \tcode{CAPABILITY\_ISSUABLE\_RING}，未注册动作不存在记录，$\chi_1$ 或 $\chi_6$ 必然失败。复合系统仍为默认拒绝，但\textbf{完全中介者是 broker，不是网关}。
\end{assumptioninner}

\begin{assumptioninner}[{同步串行化中介}]\label{ass:a2}
$\Auth$ 在系统调用边界内同步求值；\emph{内核中介循环内}的调用被宿主事件循环全序化；$\chi$ 求值与处理器调用之间不会交错撤销写（网关检查内部无 TOCTOU 窗口）。注意该全序声明只覆盖内核中介循环——执行域中的 worker 线程/子进程代码不受其约束。
\end{assumptioninner}

\begin{assumptioninner}[{无缓存}]\label{ass:a3}
九个检查在每次调用上重新求值；任何决策不跨系统调用缓存。
\end{assumptioninner}

TOCTOU（time-of-check to time-of-use）正是 A2/A3 所防御的漏洞类：名字与对象的绑定可能在检查与使用之间改变~\cite{bishop1996}。ConsistenCy 的对策是三重的：逐调用重评（A3）、检查与执行的中介点串行化（A2）、以及作用域钉定\emph{内容}（sha）而非可变名字。

\section{安全不变式 S1}

记 $\Exec(\alpha,r,t)$ 为“请求 $(\alpha,r)$ 的处理器在步 $t$ 产生效果”；每次执行事件携带其中介点 $m(q)$——同一 \tcode{gateway.invoke} 内 $\Auth$ 被求值的串行化步，$m(q)\prec_{\mathrm{ser}}t$。

\begin{proposition}[安全不变式 S1：边界授权]\label{prop:s1}
在假设 A1--A3 下，
\begin{equation}
\mathbf{S1}\;:\quad \Box\bigl(\,\Exec(\alpha,r,t)\;\Rightarrow\;\Auth(q,m(q))=1\,\bigr).
\end{equation}
\end{proposition}
\begin{proof}[证明梗概（构造性）]
在 $\mathcal{M}_C$ 中，产生 $\Exec$ 事件的唯一迁移规则是“invoke 求值 $\Auth$；为 0 则抛出 \tcode{CapabilityError}（处理器不被调用）；为 1 才进入处理器”（\tcode{syscall/authorize.ts:46--102}；intent 类在授权\emph{之前}即抛出，内联处理器调用数为零）。若 A1--A3 任一失守，该规则就不再是执行事件的唯一生产者，S1 变为空洞。
\end{proof}
\tagm{ENGINEERING\_INVARIANT（代码可证调用路径；“不存在旁路”是对全部未来代码的全称否定，属假设）}

\paragraph{失败条件。}
(i) 绕过 \tcode{invoke} 的效果路径（A1 失守；注意 OS 级文件系统/网络围堵在当前实现中除子进程隔离外\emph{未强制}，S1 对模型外效果不做声明）；(ii) 跨边界缓存（A3）；(iii) 非串行化运行时中的检查--使用间隙（A2 失守时 S1 在中介点语义下仍成立，但更强的“$t_r$ 后无执行”（S2）需要 A2）；(iv) 网关穿透本身不是 S1 的违反（broker 拒绝），但把完备性负担集中在单一组件上。

\section{撤销支配（S2）}\label{sec:revocation}

\begin{assumptioninner}[{标志单调}]\label{ass:a4}
$C_t$ 的迁移规则永不复位 $revoked$（撤销永久——代码可证，\tcode{broker.ts:195--218}）。
\end{assumptioninner}
\begin{assumptioninner}[{原子撤销写}]\label{ass:a5}
对 $C_{t}(h^\ast).revoked$ 的写相对于中介求值是原子的（事件循环串行化，与 A2 同基础设施）。
\end{assumptioninner}

\begin{proposition}[撤销支配]\label{prop:revocation}
设 $\mathrm{revoke}(h^\ast,t_r)$ 为 $revoked:\bot\to\top$ 的迁移（日志+事件总线；无级联；向 Cordis 的传播延迟到微任务队列）。则对每个引用 $h^\ast$ 且 $m(q)\succeq_{\mathrm{ser}}t_r$ 的请求 $q$：
\begin{equation}
\Auth(q,m(q))=0,\qquad\text{与 } K_{m(q)},V_{m(q)},P_{m(q)},\mathfrak{A}_{m(q)},Q_{m(q)}\text{ 无关}.
\end{equation}
\end{proposition}
\begin{proof}
$\Auth=\bigwedge_i\chi_i\le\chi_2$。由 A4，对一切 $t\succeq t_r$ 有 $revoked=\top$，故 $\chi_2=0$，$\Auth=0$（拒绝原因为 \tcode{revoked}，链上第 2 位）。无关性是\emph{语法性}的：$\Auth$ 的自由变元含于 $\{C_t,B_t,h,a,\alpha,r,s,t\}$；$K,V,P,\mathfrak{A},Q$ 不出现在公式中，而 $\Auth$ 是其声明输入的全函数，故排除分量的任何变动不改变裁决。结合 S1：不存在中介点 $\succeq t_r$ 的执行。
\end{proof}
\tagm{PROPOSITION（模型级；A4/A5 代码可证）}

\paragraph{意图门 carve-out（如实声明）。}
对 $\alpha\in\Pi_{\mathrm{intent}}$，授权发生在\textbf{意图接受时}：接受之后的撤销不抹除已接受意图；不可逆外部效果在幂等键去重与下游租约/退避/重试（租约 30\,s、超时 15\,s、最多 3 次尝试）之下继续执行。形式上 $\mathrm{Accept}(key)\Rightarrow\mathrm{Effect}$ 由接受器而非效果时刻的 $\Auth$ 支配。对称地，接受时刻同样“冻结”了 $\chi_3$（过期）与预算预留的判定。因此 S2 的正确读法是\textbf{作用于中介点}；若读作“撤销后不再发生任何来自撤销前已接受意图的不可逆效果”，则该命题对 intent 类为假。

\paragraph{在飞操作。}
中介点先于 $t_r$ 的操作会完成——撤销不是对在飞副作用的取消，只是对后续授权的取消。这是“Effect $\neq$ Rollback”公理（见第~\ref{ch:sep} 章）在时序上的体现。

\section{撤销时间线}

图~\ref{fig:timeline} 展示撤销支配与异步传播的并存：内核裁决（黑）在 $t_r$ 后立即转为拒绝；Fiber 生命周期（灰）要等到 $t_{\mathrm{unload}}$ 才完成卸载；两者之间的陈旧门面窗口 $W_{\mathrm{stale}}$ 内，任何经陈旧门面发起的系统调用都被拒绝（命题~\ref{prop:revocation}），窗口长度 $\delta_{\mathrm{prop}}\ge 0$ 甚至可以是 $\infty$ 而不影响安全性。

\begin{figure}[htbp]
\centering
\begin{adjustbox}{max width=\textwidth}
\begin{tikzpicture}[x=1cm,y=1cm,>={Stealth[length=2mm]},
  ev/.style={circle,fill=accentred!80,inner sep=1.6pt},
  ok/.style={font=\scriptsize},
  ln/.style={thick,inkgray}]
% axis
\draw[->,thick] (-0.2,0) -- (13.2,0) node[below right,font=\small] {$t$};
% kernel verdict lane
\draw[ln] (0,1.6) -- (13,1.6);
\node[font=\scriptsize\bfseries,anchor=west] at (0,2.25) {内核裁决 $\Auth(\cdot\mid h^\ast)$};
\node[ok,anchor=west] at (0.15,1.95) {\textcolor{accentblue!80}{ALLOW（$\chi_1..\chi_9$ 全过）}};
\node[ev] (tr) at (4.2,1.6) {};
\node[font=\scriptsize\bfseries,above=2pt of tr] {$t_r$：revoke};
\node[ok,anchor=west,text=accentred] at (4.45,1.95) {DENY（reason=\tcode{revoked}）$\;\forall\,m(q)\succeq t_r$};
% fiber lane
\draw[ln] (0,0.0) -- (13,0.0);
\node[font=\scriptsize\bfseries,anchor=west] at (0,0.62) {Fiber 生命周期};
\node[ok,anchor=west] at (0.15,0.32) {ACTIVE（持有门面）};
\node[ev,fill=inkgray] (tu) at (7.6,0) {};
\node[font=\scriptsize,above=2pt of tu] {$t_{\mathrm{unload}}$：dispose$\to$UNLOADING$\to$PENDING};
\node[ok,anchor=west,text=inkgray] at (7.85,0.32) {PENDING（资源回收完成）};
% stale window
\begin{scope}[on background layer]
\fill[softrose] (4.2,-0.35) rectangle (7.6,0.9);
\end{scope}
\draw[decorate,decoration={brace,amplitude=4pt},inkgray] (4.2,-0.5) -- (7.6,-0.5) node[midway,below=5pt,font=\scriptsize\itshape] {$W_{\mathrm{stale}}$：陈旧门面窗口，$\delta_{\mathrm{prop}}\ge 0$（可为 $\infty$）};
\node[font=\tiny\itshape\color{inkgray},align=center] at (5.9,0.14) {经陈旧门面的 syscall\\一律 DENY};
% denied syscalls marks
\foreach \x in {4.9,5.7,6.6,7.3} \node[ev,fill=inkgray!60] at (\x,-0.0) {};
\foreach \x/\y in {4.9/-0.8,5.7/-0.8,6.6/-0.8} \node[font=\tiny,text=accentred] at (\x,\y) {$\times$};
% accepted intent lane
\draw[ln,dashed] (0,-1.6) -- (13,-1.6);
\node[font=\scriptsize\bfseries,anchor=west] at (0,-2.25) {已接受意图（carve-out）};
\node[ok,anchor=west] at (0.15,-1.9) {intent 于 $t_a<t_r$ 被接受（$\Auth$ 在 $t_a$ 判定）};
\node[ev,fill=accentblue!80] (ta) at (2.4,-1.6) {};
\node[font=\scriptsize,above=2pt of ta] {$t_a$};
\node[ok,anchor=west,text=inkgray] at (6.5,-1.9) {效果经幂等键+租约继续执行（不受 $t_r$ 影响）};
\end{tikzpicture}
\end{adjustbox}
\caption{能力授权时间线与撤销。上：内核裁决在 $t_r$ 后同步转为拒绝（命题~\ref{prop:revocation}）；中：Fiber 卸载异步完成，窗口 $W_{\mathrm{stale}}$ 内的调用全部被拒；下：撤销前已接受的意图不受影响（carve-out）。}
\label{fig:timeline}
\end{figure}

\section{与经典理论的差异点}

与 Redell--Fabry 的选择性撤销相比，ConsistenCy 选择了\emph{更简单}的语义：撤销是全局永久的，不需要“看守人（caretaker）”间接层来隔离单个持有者；这换来的是命题~\ref{prop:revocation} 的单调性证明极其简短。与 seL4 相比，ConsistenCy 没有任何机器检查的证明——S1/S2 是“模型级命题 + 源码核验的构造性论据”，其“不存在旁路”部分（A1 的全称否定）恰恰是形式化验证（如 Isabelle/HOL 之于 seL4~\cite{klein2009}）才能消除的假设负担；这构成本报告在第~\ref{ch:limits} 章中声明的核心局限之一，也是后续模型检验扩展（\S\ref{sec:modelcheck}）的动机。

% 第 5-6 章：Cordis 生命周期/上下文可见性与授权的正交性
\chapter{Cordis Lifecycle 与 Authorization 的分离}\label{ch:sep}

\section{工程事实}

Cordis 层（Cordis 4.0.0-rc.8 库语义之上）管理响应式组件的生命周期：Fiber 状态为 PENDING / LOADING / ACTIVE / UNLOADING；共效（Coeffect）以声明式依赖（如 \tcode{Plugin.Function.inject=["ast"]}）表达“进入 ACTIVE 所需的服务”；每个 Agent 经 \tcode{root.isolate} 获得隔离的依赖上下文；内核能力经 \tcode{CapabilityLifecycleAdapter} 桥接为服务门面（facade）。关键事实是\textbf{两个时间尺度}：
\begin{itemize}[itemsep=0.05em]
  \item 能力撤销后，\tcode{capability.revoked} 事件在\emph{微任务队列}上调度销毁器；在传播完成前，Fiber 可以保持 ACTIVE 并持有\emph{陈旧门面}；
  \item 但该陈旧门面在\textbf{下一次系统调用}处即被内核拒绝（命题~\ref{prop:revocation}）。
\end{itemize}
即：撤销在生命周期层是\emph{最终一致}的，在授权层是\emph{同步}的。组件公理“安全性不依赖传播时序”由此获得精确含义。

\section{理论映射}

Cordis 层在 Harel--Pnueli 的分类下是一个\emph{响应式系统}（reactive system）——持久运行、事件驱动、与环境持续交互而非计算到终态~\cite{harel1985}；这一定位把“Fiber 生命周期为何需要独立于任何单次运行来推理”变得自然。Fiber 的监督式生命周期（依赖缺失即不激活、失败即清理、可被监督者处置）在理论上最接近两条线：其一是 Actor 模型的严格语义~\cite{agha1986} 与 Erlang/OTP 的监督树（supervision tree）——Armstrong 的论文把“错误内核 vs 错误携带”与分层重启发展为可靠分布式系统的构造原理~\cite{armstrong2003}；其二是 Miller 的 vat 模型：事件循环受限的响应式并发与访问控制的统一~\cite{miller2006}（vat 的技术核心之一是 vat 间消息传递的全序保证；ConsistenCy 的对应物是内核中介循环的串行化——假设 A2——但 Cordis 层并未声明跨上下文的端到端消息序，故本类比限定于“事件循环受限的并发单元 + 无环境固有权限”这一结构层）。\textbf{映射强度：\STRONG}。必须指出的\emph{分歧点}：OTP 的监督哲学是“让它崩溃、以重启掩盖故障”；ConsistenCy 的触发计划执行是\textbf{单次尝试 + 诚实的重启失败记录}（执行中断在启动恢复时被标记为 \tcode{failed} 而非静默重试）——这是同一监督传统内的另一个设计点，选择“失败关闭的诚实”而非“尽力掩盖的可用性”。该分歧不能被表述为“ConsistenCy 实现了 OTP 监督语义”，只能表述为“处于同一监督谱系”。

\section{可用性与授权的正交性}

\begin{definition}[可用性]\label{def:avail}
$\Avail(a,s,t)\;\equiv\;$[共效服务 $s$ 在 Agent $a$ 的隔离上下文中已提供] $\wedge$ [$a$ 的 Fiber 处于 ACTIVE]。
\end{definition}

\begin{proposition}[可用性 $\nvdash$ 授权，反之亦然]\label{prop:orthavail}
$\Avail(a,s,t)\nvdash \Auth(a,\alpha,r,s,t)$ 且 $\Auth\nvdash\Avail$。且安全性只依赖 $\Auth$：若 $\Avail\equiv 0$，则无任何执行发生，S1 空洞成立；可用性失败只降低\emph{活性}（效果无法完成、Fiber 停留 PENDING、Agent 甚至无法发起调用），永不授予\emph{权威}。
\end{proposition}
\begin{proof}[见证结构]
$\mathcal{W}_1$（授权 $\wedge$ 不可用）：$C_t$ 有有效未撤销的 (\tcode{ast.query}, snapshot) 记录，但 \tcode{ast} 服务从未提供——Fiber 停留 PENDING，Agent 永远无法调用；假想请求的 $\Auth=1$，而 $\Avail=0$。$\mathcal{W}_2$（未授权 $\wedge$ 可用）：Fiber ACTIVE 持有门面，$t_r<t$ 时底层能力被撤销而传播未运行；$\Avail=1$，$\Auth=0$（命题~\ref{prop:revocation}）。两个见证各自满足一个合取而证伪另一个，故两个蕴含都不有效。
\end{proof}
\tagm{PROPOSITION}

\begin{proposition}[传播时序无关性]\label{prop:propagation}
对一切 $\delta_{\mathrm{prop}}=t_{\mathrm{unload}}-t_r\in[0,\infty]$（含 $\infty$，即销毁器永不运行）：窗口 $W_{\mathrm{stale}}$ 内经陈旧门面发起的全部系统调用都被拒绝。$\delta_{\mathrm{prop}}$ 的上界（生命周期的有界最终一致性）只对\emph{资源回收}——一个活性/整洁性质——是必需的，对安全健全性永不需要；传播饥饿（微任务队列拥塞）打不开任何权威窗口。
\end{proposition}
\begin{proof}
窗口内每个系统调用的中介点 $\succeq t_r$；应用命题~\ref{prop:revocation}。
\end{proof}
\tagm{PROPOSITION}

\chapter{Context Visibility 与 Authority 的分离}\label{ch:ctx}

\section{工程事实}

Context VM 是“虚拟上下文内存”：不可变的 \tcode{ContextPage} 以 SHA-256 内容哈希寻址；驻留级（pinned/hot/cold/evicted）位于镜像覆盖层而非页面；工作集 $=$ pinned+hot；除 \tcode{pinned$\rightarrow$evicted} 抛错（失败关闭）外驻留级可自由迁移；COW fork 只复制页面表、共享页对象；渲染顺序由类别先序与页 id 决定（\emph{绝不}依赖哈希表插入序）。源码中明文写明：“Context VM 是状态管理，不是授权；页面成员资格不授予任何能力”（\tcode{context/manager.ts:16--19}）。当前实现\emph{没有}自动逐出、摘要或检索策略——任何暗示存在此类策略的描述都是过度陈述。

\section{可见性与权限的正交性}

\begin{definition}[可见性]\label{def:visible}
$\Vis(a,x,t)\;\equiv\;\mathrm{PT}_a(x)\in\{\text{pinned},\text{hot}\}$（页 $x$ 在 $a$ 的工作集中）。
\end{definition}

\begin{proposition}[可见性 $\nvdash$ 授权，反之亦然]\label{prop:orthvis}
$V_t$ 不是 $\Auth$ 的输入（定义~\ref{def:auth} 的自由变元不含 $V_t$）。见证：页在而无能力（“页面成员资格不授予能力”）给出 $\Vis\wedge\neg\Auth$；有 \tcode{repo.read} 能力而相关页全部被逐出到 cold 给出 $\Auth$-可能 $\wedge\neg\Vis$。
\end{proposition}
\tagm{PROPOSITION}

\begin{figure}[htbp]
\centering
\begin{adjustbox}{max width=\textwidth}
\begin{tikzpicture}[x=1cm,y=1cm]
% axes
\draw[->,thick] (0,0) -- (9.6,0) node[below right,font=\small] {$\Vis$（上下文可见性：Agent 看见什么）};
\draw[->,thick] (0,0) -- (0,6.4) node[above,font=\small,align=center] {$\Auth$（内核权威：Agent 被允许做什么）};
% quadrant labels
\node[draw=framegray,fill=softrose,rounded corners=2pt,align=center,font=\scriptsize,minimum width=3.9cm,minimum height=1.7cm] (q2) at (2.3,4.6) {\textbf{“看得见，做不了”}\\机密在上下文中但无能力\\受保护 syscall 仍被拒绝};
\node[draw=framegray,fill=softblue,rounded corners=2pt,align=center,font=\scriptsize,minimum width=3.9cm,minimum height=1.7cm] (q1) at (7.0,4.6) {\textbf{“看得见也做得了”}\\正常已授权读取\\（唯一执行区）};
\node[draw=framegray,fill=boxgray,rounded corners=2pt,align=center,font=\scriptsize,minimum width=3.9cm,minimum height=1.7cm] (q3) at (2.3,1.7) {\textbf{“看不见也做不了”}\\无上下文且无能力\\空洞安全区};
\node[draw=framegray,fill=softamber,rounded corners=2pt,align=center,font=\scriptsize,minimum width=3.9cm,minimum height=1.7cm] (q4) at (7.0,1.7) {\textbf{“看不见，但有权”}\\页被逐出仍可授权调用\\（pageIn 恢复后执行）};
% annotations
\node[font=\tiny\itshape\color{inkgray},align=left,anchor=west] at (0.2,-1.35) {安全只由纵轴决定（S1 只读 $\Auth$）；可见性只影响“调用是否可能被发起/上下文是否在手”。};
\node[font=\tiny\itshape\color{inkgray},anchor=west] at (0.2,-0.75) {同理 $\Avail$（服务可用性）与 $\Auth$ 正交（命题~\ref{prop:orthavail}）：可用性只影响活性。};
\end{tikzpicture}
\end{adjustbox}
\caption{上下文可见性与权威的正交轴。Context VM $\neq$ 权限系统：两轴独立变化，安全性只由纵轴（$\Auth$）承担；把横轴当安全边界的任何论述（例如“Agent 没看到密钥所以安全”）在模型上不成立。}
\label{fig:orthogonal}
\end{figure}

\section{为什么这不是语义游戏}

图~\ref{fig:orthogonal} 的正交结构排除了两类常见的错误论证：
\begin{enumerate}[itemsep=0.1em]
  \item \textbf{“看不到即安全”}：即使某秘密页从未进入 Agent 工作集，也不能据此声称机密性——那需要信息流不干涉类的性质，而本系统并未声明（密钥的机密性由服务端持有与脱敏保障，是另一条防线，不属于 $\Vis/\Auth$ 轴）；
  \item \textbf{“看到了就危险”}：页在上下文中\emph{不}提升任何权威——受保护系统调用照样要过 9 步链。反过来，$\Vis$ 的驱逐/钉定只影响性能与正确性（如 \tcode{pinned$\rightarrow$evicted} 失败关闭保护关键上下文不丢失），不影响权限。
\end{enumerate}
这与 Saltzer--Schroeder 的“关注点分离/机制经济性”一脉相承~\cite{saltzer1975}：可见性（内存语义）与权威（保护语义）是两个正交的轴，各自的最优设计互不干扰。值得注意的是，COW fork 的“父修改不流入既有子镜像”与内容寻址一起，使 $\Vis$ 成为一个\emph{纯函数式}的投影——渲染是 (页表, 驻留级, 先序) 的确定性函数，这为第~\ref{ch:evidence} 章的可复现性论证提供了上下文侧的对应物。

% 第 7 章：Agent / Run / Fiber 状态模型
\chapter{Agent / Run / Fiber 状态模型}\label{ch:states}

\section{三个状态机，三种所有者}

ConsistenCy 的运行时由三个显式状态机组成，分属不同所有者：
\begin{enumerate}[itemsep=0.1em]
  \item \textbf{ACB（Agent Control Block）LTS}：12 状态，由 KernelScheduler \emph{独占写入}。状态集 $\States$ 含 12 个状态：NEW、READY、RUNNING、五个等待态（WAIT\_LLM / WAIT\_TOOL / WAIT\_IO / WAIT\_AGENT / WAIT\_HUMAN）、SUSPENDED 与三个终态；终态集 $T_{\mathrm{term}}=\{\tcode{SUCCEEDED},\tcode{FAILED},\tcode{CANCELLED}\}$ 吸收（迁移表 \tcode{agent/state.ts:27--50}）。执行器/Agent 代码只能\emph{请求}阻塞与唤醒（\tcode{RUNNING}$\to$\tcode{WAIT}$^\tau$ 与 \tcode{WAIT}$^\tau\to$\tcode{READY} 的原因），且必须经调度器原语生效；网关\emph{不拥有}任何 ACB 迁移（其拒绝结果作为执行器层的失败请求传播）。
  \item \textbf{Run LTS}：\tcode{CREATED}$\to$\tcode{ACTIVE}$\to$终态；运行取消级联到所有非终态 Agent；\emph{已派发的外部工作不被回滚}。
  \item \textbf{Fiber LTS}：4 状态 PENDING / LOADING / ACTIVE / UNLOADING，语义由 Cordis 库拥有（外部公理；升级 Cordis 版本会重新打开本章结论）。
\end{enumerate}

理论上，ACB 的 12 状态迁移表就是 Keller 意义下的\textbf{带标迁移系统（labelled transition system, LTS）}~\cite{keller1976}；对迁移不变式（“没有任何迁移离开 \tcode{SUCCEEDED}”、“\tcode{WAIT}$^\tau$ 最终恢复或失败”）的自然语言是 Pnueli 的\textbf{时序逻辑}~\cite{pnueli1977}。映射强度：\DIRECT{}（不是类比——它字面上就是有限 LTS）。

\section{单向耦合积}

\begin{definition}[单向耦合 $\varphi$]\label{def:coupling}
部分函数 $\varphi$ 把调度器标签映射为 Fiber 标签：$\varphi(\mathrm{admit})=\mathrm{activate}$，$\varphi(\mathrm{suspend})=\mathrm{cleanup}$，$\varphi(\mathrm{cancel})=\mathrm{dispose}$；内核事件 $\mathrm{revoke}$ 自主触发 $\mathrm{dispose}$。方向性：内核事件可以\emph{禁用} Fiber 迁移（撤销禁用后续激活并强制销毁），但\emph{永不启用}任何调度器未准入的激活；没有任何 Fiber 事件改变 ACB 状态。
\end{definition}

\begin{proposition}[积的良好定义与派生不变式]\label{prop:product}
$\varphi$-受限异步积 $\mathcal{T}=\mathcal{T}_{\mathrm{ACB}}\parallel_\varphi\mathcal{T}_{\mathrm{F}}$（两迁移关系的交错，Fiber 激活/清理标签须与其 $\varphi$-原像共现，dispose 允许自主发生）是良定义 LTS：$\varphi$ 是不相交标签集之间的部分函数，复合迁移关系在每个复合标签处单值。状态空间 $\subseteq\States\times S_{\mathrm{F}}$（48 对）。由桥接映射派生：$\text{Fiber ACTIVE}\Rightarrow \text{ACB}\in\{\tcode{RUNNING}\}\cup\{\tcode{WAIT}^\tau\}$；故 $(\tcode{WAIT\_LLM},\text{ACTIVE})$ 可达（已证事实），而 $(\tcode{NEW},\text{ACTIVE})$ 与 $(\tcode{SUSPENDED},\text{ACTIVE})$ 不可达。
\end{proposition}
\begin{proof}
$\varphi$ 的部分性与函数性保证合成约束无歧义；不可达性由 $\varphi$-共现要求直接得出（\tcode{NEW} 与 \tcode{SUSPENDED} 不出现在任何 $\varphi$ 定义域原像的使能条件中）。
\end{proof}
\tagm{PROPOSITION（LTS 定义为 ENGINEERING\_INVARIANT；积性质为模型级）}

$(\tcode{WAIT\_LLM},\text{ACTIVE})$ 的可达性值得强调：它说明 \emph{Fiber ACTIVE $\neq$ 调度器 RUNNING} 是系统的\emph{设计事实}而非缺陷——Fiber 反映 harness 侧的组件存活，调度器跟踪协作式等待。这一分离与第~\ref{ch:sep} 章的正交性共同构成“安全看内核、活性看生命周期”的二分。

\begin{figure}[htbp]
\centering
\begin{adjustbox}{max width=\textwidth}
\begin{tikzpicture}[
  st/.style={draw=framegray,fill=boxgray,rounded corners=2pt,font=\scriptsize,minimum height=1.5em,inner sep=3.5pt},
  tr/.style={draw=accentred!70,fill=softrose,rounded corners=2pt,font=\scriptsize\bfseries,minimum height=1.5em,inner sep=3.5pt},
  fs/.style={draw=framegray,fill=softgreen,rounded corners=2pt,font=\scriptsize,minimum height=1.5em,inner sep=3.5pt},
  ln/.style={-{Stealth[length=1.8mm]},inkgray!80},
  lb/.style={font=\tiny\itshape\color{inkgray}, midway, fill=white, inner sep=1pt}]
% ---- ACB LTS (left) ----
\node[font=\small\bfseries] at (3.1,4.6) {ACB LTS（12 状态，调度器所有）};
\node[st] (new) at (0.4,3.7) {NEW};
\node[st] (rdy) at (2.2,3.7) {READY};
\node[st] (run) at (4.2,3.7) {RUNNING};
\node[st] (sus) at (0.4,2.2) {SUSPENDED};
\node[tr,minimum width=3.2cm] (wait) at (3.1,2.2) {WAIT\_LLM / TOOL / IO / AGENT / HUMAN};
\node[tr] (ok) at (5.8,3.7) {SUCCEEDED};
\node[tr] (fail) at (5.8,2.2) {FAILED};
\node[tr] (canc) at (3.1,0.9) {CANCELLED};
\draw[ln] (new) -- (rdy) node[lb,above]{register};
\draw[ln] (rdy) -- (run) node[lb,above]{admit};
\draw[ln] (run) -- (wait) node[lb,right=1pt]{wait};
\draw[ln] (wait) -- (rdy) node[lb,left=1pt]{wake};
\draw[ln] (sus) -- (2.0,3.05) node[lb,left=1pt]{resume};
\draw[ln] (rdy.south west) -- (sus.north east) node[lb,left=1pt]{suspend};
\draw[ln] (run) -- (ok) node[lb,above]{succeed};
\draw[ln] (wait) -- (fail) node[lb,below]{fail};
\draw[ln] (run.south) to[bend left=8] (canc.east) node[lb,right]{cancel};
\draw[ln] (rdy.south) to[bend right=8] (canc.west) node[lb,left]{cancel/deadline};
\node[font=\tiny\itshape\color{inkgray},align=center] at (3.1,0.15) {终态吸收；截止期到期的 READY 在 ACTIVE 运行中被 CANCEL（P1 路pass），绝不准入};
% ---- Fiber LTS (right) ----
\node[font=\small\bfseries] at (10.6,4.6) {Fiber LTS（4 状态，Cordis 库所有）};
\node[fs] (pend) at (9.2,3.4) {PENDING};
\node[fs] (load) at (11.2,3.4) {LOADING};
\node[fs] (act) at (11.2,2.0) {ACTIVE};
\node[fs] (unl) at (9.2,2.0) {UNLOADING};
\draw[ln] (pend) -- (load) node[lb,above]{facade 提供};
\draw[ln] (load) -- (act) node[lb,right=1pt]{};
\draw[ln] (act) -- (unl) node[lb,below]{dispose / revoke};
\draw[ln] (unl) -- (pend) node[lb,left=1pt]{cleanup};
\draw[ln,dashed,inkgray!60] (act.west) to[bend left=28] (pend.north west) node[lb,above left=1pt]{依赖从未提供};
\node[font=\tiny\itshape\color{inkgray},align=center] at (10.6,0.9) {共效缺失 $\Rightarrow$ PENDING 吸收；\\撤销 $\Rightarrow$ dispose（异步，经微任务）};
% coupling arrows
\draw[ln,accentred!70,dashed,thick] (run.north east) to[bend left=16] node[lb,above,font=\tiny\color{accentred}]{ $\varphi(\mathrm{admit}){=}\mathrm{activate}$} (act.north west);
\draw[ln,accentred!70,dashed,thick] (sus.north west) to[bend left=40] node[lb,below left,font=\tiny\color{accentred}]{$\varphi(\mathrm{suspend}){=}\mathrm{cleanup}$} (unl.south west);
\end{tikzpicture}
\end{adjustbox}
\caption{ACB/Fiber 状态迁移图。左：12 状态 ACB（红框为终态）；右：4 状态 Fiber（绿色）。红色虚线为单向耦合 $\varphi$：内核事件可以禁用 Fiber 迁移、强制销毁，但永不启用未被调度器准入的激活；没有 Fiber 事件改变 ACB 状态（命题~\ref{prop:product}）。}
\label{fig:states}
\end{figure}

\section{模型检验扩展（\PROPOSED）}\label{sec:modelcheck}

48 对的有限积空间提示了一个尚未接线的验证路径：把 S1--S3 与积 LTS 的不变式编码为 CTL 公式并模型检验——这正是 Clarke--Emerson 分支时序逻辑模型检验的诞生场景~\cite{clarke1981}。当前实现以测试钉住若干不变式（终态吸收、重复入队防御、不可达组合），但未做穷举检验。该扩展不改变任何结论的效力，只是把“逐条断言”升级为“穷举排除”；列入第~\ref{ch:future} 章的研究议程。

% 第 8 章：Workflow DAG 与执行语义
\chapter{Workflow DAG 与执行语义}\label{ch:workflow}

\section{四段管线}

工作流运行时把一个经验证的定义变成一次受治理的运行，管线为四段（全部源码核验）：
\begin{enumerate}[itemsep=0.1em]
  \item $\operatorname{validate}:\mathrm{Def}\to\mathrm{Def}\cup\{\bot\}$——结构良构性校验，含以 Kahn 入度消解实现的环检测；成功时给出确定性拓扑序 $\sigma$（同层平局按节点 id 决断）；
  \item $\operatorname{compile}:\mathrm{Def}\times\mathcal{R}\to\mathrm{Plan}\cup\{\bot\}$——节点类型对运行时注册表 $\mathcal{R}$ 解析（当前两类：\tcode{analyzer.deterministic-evidence} 需 $\{\tcode{repo.read},\tcode{evidence.write}\}$；\tcode{verifier.persisted-evidence} 需 $\{\tcode{evidence.read}\}$）；能力需求必须是已注册系统调用动作；\textbf{编译器不发放任何持久授权}——需求是数据；
  \item $\operatorname{dryload}$——逐节点可行性复查（复用同一编译函数；明示“仅可行性检查”）；
  \item $\operatorname{execute}$——\tcode{host.trigger} 在\emph{任何}运行记录之前失败关闭（未知仓库 404 / 快照不可得 503 / 定义不可执行 409）；钉定 HEAD 后创建快照；持久化运行后异步执行；节点内\textbf{逐系统调用授权}（第~\ref{ch:security} 章）。
\end{enumerate}
定义携带字面量 \tcode{failurePolicy="fail-closed"}：节点失败产生终态运行状态，而非部分继续。

\section{拓扑序定理与 WF-net 负结果}

\begin{theorem}[拓扑序存在性]\label{thm:topo}
无环有向图 $G=(V,E)$，$|V|\ge1$，存在拓扑序 $\sigma:V\to\{1..|V|\}$ 使 $(u,v)\in E\Rightarrow \sigma(u)<\sigma(v)$；Kahn 算法以 $O(|V|+|E|)$ 时间计算一个拓扑序（或检出环）。
\end{theorem}
\tagm{THEOREM\_FROM\_KNOWN\_MODEL（算法课标准结论；实现即 \tcode{validate.ts} 的入度消解）}

\paragraph{Petri 网/WF-net 负结果（\NA-for-now）。}
工作流理论的标准形式化是 van der Aalst 的工作流网（WF-net）及其\emph{健全性（soundness）}：正常完成可达、无死任务、无遗留托肯~\cite{vanderaalst1998}。经评估，对\textbf{当前的} ConsistenCy 语义——无环 DAG、确定性平局决断、无延迟选择（deferred choice）与里程碑——Kahn 拓扑序已经给出终止性与无死锁的构造性保证，WF-net 的位置不变量与托肯演算\textbf{不增加证明力}。该映射记为负结果：完整 Petri 机制\emph{推迟}，仅当语言将来引入环、迭代或延迟选择时才值得引入。这不是对 WF-net 理论的否定，而是对其适用边界的诚实划定。

\begin{figure}[htbp]
\centering
\begin{adjustbox}{max width=\textwidth}
\begin{tikzpicture}[
  ph/.style={draw=accentblue!70,fill=softblue,rounded corners=3pt,font=\scriptsize,align=center,minimum height=2.2em,inner sep=5pt,minimum width=2.35cm},
  ac/.style={draw=framegray,fill=boxgray,rounded corners=2pt,font=\scriptsize,align=center,minimum height=2.2em,inner sep=4pt},
  ln/.style={-{Stealth[length=2mm]},inkgray!80,thick},
  lb/.style={font=\tiny\itshape\color{inkgray},midway,fill=white,inner sep=1pt}]
\node[ph] (def) at (0,2.3) {Definition\\（数据，无权威）};
\node[ph] (val) at (3.3,2.3) {validate\\Kahn 环检测\\拓扑序 $\sigma$};
\node[ph] (cmp) at (6.6,2.3) {compile\\注册表解析\\能力需求=数据};
\node[ph] (dry) at (9.9,2.3) {dry-load\\可行性（仅检查）};
\node[ac,fill=softgreen!50] (trig) at (0,0) {触发源\\manual / repository\_change};
\node[ac,fill=softgreen!50] (plan) at (3.3,0) {持久触发计划\\哈希去重键 + 围栏单飞};
\node[ac] (host) at (6.6,0) {host.trigger\\失败关闭前置检查};
\node[ac,fill=softamber!60] (run) at (9.9,0) {Run（运行）\\状态 running$\to$终态};
\node[ac,fill=softrose] (gate) at (9.9,-2.0) {逐 syscall 的 $\Auth$\\（唯一权威点）};
\node[ac,fill=softamber!60] (ev2) at (13.2,0) {Evidence / Finding\\（接地后落地）};
\draw[ln] (def) -- (val) node[lb,above]{良构};
\draw[ln] (val) -- (cmp) node[lb,above]{$\sigma$};
\draw[ln] (cmp) -- (dry);
\draw[ln] (trig) -- (plan) node[lb,above]{eventId=观测};
\draw[ln] (plan) -- (host);
\draw[ln] (host) -- (run);
\draw[ln] (dry) -- (9.9,1.0) node[lb,right=2pt]{可行性通过};
\draw[ln] (run) -- (gate) node[lb,right]{每个受保护操作};
\draw[ln] (gate) -- (ev2) node[lb,right]{允许才落地};
\draw[ln,dashed,inkgray!50] (def.south) to[bend right=15] node[lb,below]{数据 $\not\Rightarrow$ 裁决} (gate.west);
\end{tikzpicture}
\end{adjustbox}
\caption{工作流 DAG $\rightarrow$ 校验 $\rightarrow$ 运行时管线。上排（蓝）是纯数据变换；下排从触发到运行进入真实世界，但唯一的权威点是每个系统调用处的 $\Auth$。虚线标注被排除的推理路径：定义层命题不能蕴含任何授权结果。}
\label{fig:workflow}
\end{figure}

\section{裁决面独立性（“校验 $\neq$ 授权”）}

\begin{proposition}[裁决面独立性（单步、固定状态）]\label{prop:verdictplane}
把状态分解为 $X_t=(\mathcal{D}_t,Y_t)$，其中 $\mathcal{D}_t$ 是数据面坐标（定义、编译计划、绑定、触发模式、事件标识），$Y_t\supseteq(C_t,B_t)$。裁决泛函 $\mathcal{V}:(q,X)\mapsto\Auth(q)$ 在数据替换下不变：对一切 $\mathcal{D},\mathcal{D}'$，
\begin{equation}
\mathcal{V}\bigl(q,(\mathcal{D},Y)\bigr)=\mathcal{V}\bigl(q,(\mathcal{D}',Y)\bigr).
\end{equation}
因此，任何局限于\emph{同一时刻}求值的数据面命题 $\varphi$ 都不能蕴含任何非平凡的授权结果（除非 $\varphi$ 本身不可满足或恒真）。
\end{proposition}
\begin{proof}
由定义~\ref{def:auth} 的语法：$\Auth$ 的自由变元不含任何 $\mathcal{D}$-坐标。见证：(i) 合法计划 + 零能力发放 $\Rightarrow$ 运行内所有请求被拒（$\chi_1$ 失败）；(ii) 无计划（手动运行）+ 内核发放的能力 $\Rightarrow$ 请求被允许。校验对任何授权结果既不必要也不充分。
\end{proof}
\tagm{PROPOSITION（安全不变式 S3 的形式内容）}

\paragraph{适用范围（对抗误读）。}
本命题\textbf{不是}轨迹层的不干涉（noninterference）：数据面的授权标志（catalog 投射到 \tcode{ReviewWorkload.issueCapabilities}）\emph{因果地决定}哪些能力被申请发放，从而影响\emph{后续}裁决的流。声明的、且仅声明的是：发放本身是内核侧经 Ring 策略校验的操作（Ring 0 动作不能在内核外发放），因此“数据 $\to$ 权威”的唯一合法通道是\textbf{内核侧的 Ring 校验发放}；对任何\emph{给定}请求的裁决不随数据面直接变化。触发事件是“纯观测溯源，绝非授权”（\tcode{triggers.ts}）。

\section{执行健全性与活性}

\begin{proposition}[执行健全性与活性]\label{prop:execsound}
\textbf{健全性}：运行的每个外部效果要么 (a) 是中介点处 $\Auth=1$ 的中介系统调用（S1 在工作流层的实例化），要么 (b) 是经幂等、哈希承诺、持久 Sink 落地的已接受提交意图。\textbf{活性}：在 (F1) $G$ 有限无环；(F2) 节点执行有限步内终止（确定性分析器全函数；LLM 路径一次修复尝试后类型化失败）；(F3) fail-closed 使节点失败产生终态；(F4) 公平派发（\emph{期望性}假设：单飞、5\,s 轮询、批量 5、单次尝试、无重试；高优先级干扰可饿死）；(F5) 无崩溃（启动恢复把 \tcode{executing} 标记为 \tcode{failed}——诚实失败而非进展）之下，运行按 $\sigma$ 序在有限步内到达终态。
\end{proposition}
\tagm{健全性 ENGINEERING\_INVARIANT（逐路径代码可证）；活性 PROPOSITION（F4/F5 显式为期望性假设；F2 的提供商挂起已并入 $P_t$ 可用性假设）}

\paragraph{诚实限制（如实编码的实现边界）。}
\begin{itemize}[itemsep=0.05em]
  \item \textbf{触发是“恰好一次发起、至多一次执行”}：持久去重键 $\mathrm{SHA256}[\text{domain},\text{repo},\text{def},\text{event}]$ 与围栏式 \tcode{UPDATE WHERE status='pending'} 保证每个事件每个定义至多发起一次运行；但计划执行\emph{单次尝试、无重试循环}——鲁棒性以“状态转换间崩溃”为界，不含运行中挂起。
  \item \textbf{CommitCoordinator 的去重态在内存}：未配置持久 Sink 的重启会丢失去重记忆，重启后重复接受同一幂等键成为可能（“恰好一次发起”以 Sink 持久性为条件）。
  \item \textbf{监督侧事件去重也在内存}：事件\emph{标识}仅经派生键持久。
\end{itemize}
这三条都不是缺陷声明，而是把“什么在什么故障模型下成立”写清楚的边界条件；对应的研究问题（恰好一次效果发起的可用形式化）见第~\ref{ch:future} 章。

% 第 9 章：Queueing / Scheduling / Resource Governance
\chapter{排队、调度与资源治理}\label{ch:queueing}

\section{调度器的真实形状}

KernelScheduler 的准入是三路串行判定（源码核验，\tcode{scheduler.ts:225--271}）：
\begin{enumerate}[itemsep=0.05em]
  \item \textbf{P1 截止期执法}：ACTIVE 运行中已过期（$now\ge deadline$）的 READY Agent 被\emph{取消}而非准入；
  \item \textbf{P2 容量守卫}：$runningCount\ge c$（$c=$ \tcode{maxRunningAgents}，当前工作流执行器接线为 $c=1$）则不准入；
  \item \textbf{P3 优先级-先来先服务}：最大 \tcode{priority} 者胜出，严格 $>$ 平局决断 $=$ 同优先级内 FIFO——确定性准入。
\end{enumerate}
远程 LLM 推理\textbf{不可物理抢占}：调度器的真实模型是“准入控制 $+$ 协作调度 $+$ 预算/并发控制”，而不是抢占式 CPU 调度器。这一形状声明直接决定了本章理论的适用边界。

\begin{figure}[htbp]
\centering
\begin{adjustbox}{max width=\textwidth}
\begin{tikzpicture}[
  b/.style={draw=framegray,fill=boxgray,rounded corners=2pt,font=\scriptsize,align=center,minimum height=2.0em,inner sep=4pt},
  q/.style={draw=accentblue!70,fill=softblue,rounded corners=2pt,font=\scriptsize,align=center,minimum height=2.0em,inner sep=4pt},
  r/.style={draw=framegray,fill=softamber,rounded corners=2pt,font=\scriptsize,align=center,minimum height=2.0em,inner sep=4pt},
  ln/.style={-{Stealth[length=2mm]},inkgray!80,thick},
  lb/.style={font=\tiny\itshape\color{inkgray},midway,fill=white,inner sep=1pt}]
\node[b] (repo) at (0,0) {仓库\\（外部状态 $R_t$）};
\node[b] (sup) at (2.6,0) {轮询监督\\digest 去重};
\node[q] (ev) at (5.2,0) {RepositoryEvent\\（去重后）};
\node[q] (qq) at (7.9,0) {触发计划 $\to$ 运行队列 $Q_t$};
\node[q] (sch) at (10.6,0) {3 路准入\\P1/P2/P3};
\node[r] (run) at (10.6,-1.9) {Run（$c$ 并发）\\ACB/Fiber};
\node[r] (llm) at (7.9,-1.9) {LLM 网关\\预算反压};
\node[r] (evid) at (5.2,-1.9) {证据/发现\\$E_t,F_t$};
\node[b] (out) at (2.3,-1.9) {人类决策\\（WAIT\_HUMAN）};
\draw[ln] (repo) -- (sup) node[lb,above]{$\lambda_c$（变更流）};
\draw[ln] (sup) -- (ev);
\draw[ln] (ev) -- (qq);
\draw[ln] (qq) -- (sch);
\draw[ln] (sch) -- (run) node[lb,right]{$u_t=\pi$};
\draw[ln] (run) -- (llm) node[lb,above]{逐调用授权};
\draw[ln] (llm) -- (evid) node[lb,above]{接地};
\draw[ln] (evid) -- (out);
\draw[ln,dashed,inkgray!50] (sch.south) to[bend right=20] node[lb,below]{截止期取消} (out.east);
\end{tikzpicture}
\end{adjustbox}
\caption{仓库事件 $\rightarrow$ 队列 $\rightarrow$ 调度器 $\rightarrow$ 运行 $\rightarrow$ 证据的端到端管线（本章的排队对象即此管线的 $Q_t$ 与 Run 段）。}
\label{fig:pipeline}
\end{figure}

\section{排队基线模型（\textbf{基线近似}）}

本节全部内容是\textbf{显式声明的简化}：以 Kendall 记号~\cite{kendall1953} 的词汇做容量级推理。设 $\lambda$ 为评审任务到达率、$\mu$ 为单服务器完成率、$\tau=1/\mu$、$A_e=\lambda/\mu$、$\rho=\lambda/(c\mu)$。

\begin{theorem}[稳定性（$M/M/c$ 理想化的定理）]\label{thm:stability}
$\lambda<c\mu$ 是 $M/M/c$ 正常返的充要条件~\cite{kleinrock1975}。
\end{theorem}

\begin{theorem}[Erlang-C 等待（$M/M/c$ 理想化的定理）]\label{thm:erlangc}
到达需等待的概率与均值等待：
\begin{equation}
P_{\mathrm{wait}}=\frac{\dfrac{A_e^{c}}{c!\,(1-\rho)}}{\displaystyle\sum_{k=0}^{c-1}\frac{A_e^{k}}{k!}+\dfrac{A_e^{c}}{c!\,(1-\rho)}},\qquad
\E[W_q]=\frac{P_{\mathrm{wait}}}{c\mu-\lambda},\qquad \E[W]=\E[W_q]+\tau.
\end{equation}
$c=1$ 时退化为 $P_{\mathrm{wait}}=\rho$、$\E[W_q]=\rho/(\mu-\lambda)$~\cite{kleinrock1975}。
\end{theorem}

\begin{theorem}[Little 定律]\label{thm:little}
$L=\lambda W$ 与 $L_q=\lambda W_q$：平均在队数 $=$ 到达率 $\times$ 平均逗留时间，\emph{与排队纪律无关}（优先级下同样成立）~\cite{little1961}。对 READY 积压的工程含义：测得 $L$ 与 $\lambda$ 即可估计 $W$，无需解整个系统。
\end{theorem}

\begin{theorem}[Kingman 重交通近似（GI/G/1 理想化的定理）]\label{thm:kingman}
单服务器情形（$c=1$；\textbf{非}等式，且不直接适用于 $c>1$）：
\begin{equation}
W_q\;\approx\;\frac{\rho}{1-\rho}\cdot\frac{c_a^{2}+c_s^{2}}{2}\cdot\tau,
\end{equation}
$c_a^2,c_s^2$ 为到达间隔与服务时间的平方变异系数；$M/M/1$（$c_a^2=c_s^2=1$）时精确~\cite{kingman1962}。信息：接近容量时延迟随\emph{变异性}成倍膨胀。
\end{theorem}

\begin{theorem}[非抢占 HOL 优先级（$M/G/1$ 类理想化的定理）]\label{thm:hol}
优先级类 $1>2>\cdots$，类内 FIFO（即 P3 的形状）。记 $R_k=\sum_{i\le k}\rho_i$、$W_0=\sum_i\lambda_i\,\E[S_i^2]/2$，则
\begin{equation}
\E\bigl[W_q^{(k)}\bigr]=\frac{W_0}{(1-R_{k-1})(1-R_k)}
\end{equation}
（Kleinrock 非抢占优先级均值等待~\cite{kleinrock1975}。）
\textbf{饥饿形状}：当高类饱和 $R_{k-1}\to1$ 时 $\E[W_q^{(k)}]\to\infty$——即使总体 $\rho<1$；低优先级均值延迟以 $(1-R_{k-1})^{-1}$ 爆炸。公平份额替代（逐流隔离、加权份额）在文献中成熟~\cite{demers1989}，但\textbf{未实现}——ConsistenCy 以确定性严格优先级换取简单性，把低类饥饿留给截止期取消兜底。
\end{theorem}

\begin{figure}[htbp]
\centering
\begin{adjustbox}{max width=\textwidth}
\begin{tikzpicture}
\begin{axis}[
  width=0.94\textwidth, height=6.4cm,
  xlabel={利用率 $\rho=\lambda/(c\mu)$},
  ylabel={$\E[W_q]/\tau$（无量纲）},
  xmin=0, xmax=0.97, ymin=0, ymax=14,
  legend style={font=\scriptsize,at={(0.03,0.97)},anchor=north west,draw=none,fill=none},
  grid=major, grid style={dashed,inkgray!25},
  tick label style={font=\scriptsize}, label style={font=\small}]
\addplot[very thick, accentblue, domain=0:0.95,samples=200] {x/(1-x)};
\addlegendentry{$M/M/1$ 精确：$\rho/(1-\rho)$（$c_a^2=c_s^2=1$）}
\addplot[thick, accentred!85, dashed, domain=0:0.93,samples=200] {x/(1-x)*2.5};
\addlegendentry{Kingman 变异性项：批化到达 + 重尾服务（$(c_a^2{+}c_s^2)/2=2.5$）}
\addplot[thick, inkgray!80, dotted, domain=0:0.90,samples=200] {x/(1-x)*0.6};
\addlegendentry{低变异性服务（$=0.6$，供对照）}
\end{axis}
\end{tikzpicture}
\end{adjustbox}
\caption{排队模型的\textbf{理论}曲线（theoretical queueing illustration, \emph{not} a measured ConsistenCy benchmark）：接近容量时延迟爆炸，且被服务时间变异性乘性放大（Kingman）。LLM 服务时间的实证形状（令牌量化、共批耦合、历史相关）见 \S\ref{sec:servicetime}。}
\label{fig:queueing}
\end{figure}

\paragraph{图~\ref{fig:queueing} 的系特异性推论（使其非装饰性）。}
排队曲线本身是教科书结论；把它接入 ConsistenCy 的实际架构可以得到两条可行动的推论。其一，\textbf{到达结构偏离泊松是架构决定的，容量余量必须显式计入 $c_a^2$}：监督管线对变更流做了三重变换——轮询栅格量化检测、防抖冲刷（检测时刻 $+\,T_{\mathrm{deb}}$，离栅格连续变化）、突发合并（去重把一次观测窗口内的多次变更折为至多一个待执行计划）。前两者在单仓库稀疏情形反而\emph{降低}到达间隔的变异（趋于规则化），第三者则在高峰期同时降低有效到达率与增大到达的间歇性——三个效应方向不一，任何单向断言（如“批化必然使 $c_a^2>1$”）都不成立；正确结论是\emph{$c_a^2$ 是由 $T_{\mathrm{poll}},T_{\mathrm{deb}}$ 与合并策略共同决定的、可测的量}，容量规划不应假设它为 1。其二，\textbf{Kingman 项给出参数化的余量规则}：若目标是 $\E[W_q]\le\kappa\cdot\tau$（平均等待不超过 $\kappa$ 个服务时间），由定理~\ref{thm:kingman} 解 $\frac{\rho}{1-\rho}\cdot\frac{c_a^2+c_s^2}{2}\le\kappa$，记 $S=c_a^2+c_s^2$，得
\begin{equation}
\rho\;\le\;\frac{2\kappa}{S+2\kappa},
\end{equation}
即容许利用率随 $S$ 递减——这是从该图直接读出的、以 $(\kappa,S)$ 为参数的容量上界规则：给定实测的 $S$，它能立即回答“$c=1$ 的接线最多安全承载多大的 $\lambda/\mu$”。

\subsection{假设表（含“违背”列）}

\begin{longtable}{@{}>{\raggedright\arraybackslash}p{0.24\textwidth}>{\raggedright\arraybackslash}p{0.68\textwidth}@{}}
\toprule
\textbf{假设} & \textbf{在 ConsistenCy 中的状态}\\
\midrule
\endhead
Poisson 到达 & \textbf{违背}。到达是轮询批化的：事件在轮询栅格时刻发射（每个 (repo, head, 配置摘要) 每次观测至多一个），再经防抖冲刷与去重削减；多仓库叠加\emph{趋向} Poisson，但少仓库时栅格结构主导。\\
指数服务 & \textbf{违背}。LLM 服务时间被输出令牌数量化（块状）、随共批负载耦合（非独立）~\cite{yu2022orca}、随服务器内部状态漂移（客户端不可见）~\cite{kwon2023vllm}。既非独立同分布也非指数。\\
类内 FCFS & \textbf{成立}。P3 严格平局 FIFO，代码可证的确定性。\\
单服务器“一次一个在服务” & \textbf{部分违背}。\tcode{RUNNING}$\to$\tcode{WAIT}$^\tau$ 记录 \tcode{PendingOperation} 后可释放运行位，使 LLM 调用在飞期间其他 Agent 的服务可与其重叠——$M/M/1$ 抽象在 LLM 等待窗口内失真。\\
无限缓冲 & 近似（未观察到队列容量上限）。\\
平稳性 & 存疑（仓库变更突发、负载日周期）。\\
\bottomrule
\caption{排队基线的假设状态表。整体标签：\textbf{基线近似}——诚实的记号是“带批化到达的 G/G/c 样系统”~\cite{kendall1953}；本节任何数值都不构成架构不变式。}
\label{tab:qassump}
\end{longtable}

\section{LLM 服务时间为何必须非指数建模}\label{sec:servicetime}

三个来自 LLM 服务系统的实证事实使得“把远程推理当 $M$ 服务”不可辩护：(i) Orca 的\emph{迭代级}调度与连续组批使一个请求的逗留时间取决于同批负载~\cite{yu2022orca}；(ii) PagedAttention 表明服务器端内存管理可同延迟下改变 2--4 倍吞吐——有效服务率 $\mu$ 对客户端是隐变量~\cite{kwon2023vllm}；(iii) 本报告文献核验中确认的进一步证据（KV 缓存复用使服务时间历史相关）强化同一结论。这些事实\emph{不}否定 Little 定律（它对任意纪律/分布成立），但否定了把 Erlang-C 当作预测器使用；它们正是 Kingman 变异性项（定理~\ref{thm:kingman}）值得引用的原因：预算式反压（reserve 阶段）在\emph{接近}容量之前介入，是对“$\mu$ 不可知”的正确工程响应，而不是速率限制器的缺席。

\section{CMDP 研究模型（\PROPOSED，未实现）}

把“在 token/成本/截止期约束下调度多审查 Agent”表述为约束马尔可夫决策过程~\cite{altman1999}：$\mathrm{CMDP}=(\mathcal{S},\mathcal{A},P,c,\{g_k\},\gamma)$，状态含队列构成、剩余预算、提供商可用性、运行计数；动作是准入决定（准入/推迟/取消）；$P$ 在实践中\emph{未知}（需估计，且 \S\ref{sec:servicetime} 的违背使估计困难）；目标
\begin{equation}
\min_\pi \E_\pi\Bigl[\textstyle\sum_t\gamma^t c(s_t,a_t)\Bigr]\quad\text{s.t.}\quad \E_\pi\Bigl[\textstyle\sum_t\gamma^t g_k(s_t,a_t)\Bigr]\le B_k\;\;\forall k,
\end{equation}
在有限状态/紧致动作集与 Slater 型严格可行条件下强对偶成立，鞍点处存在平稳随机最优策略~\cite{altman1999}。

\paragraph{强制诚实（三条）。}
(i) 当前调度器\textbf{不是}优化器：3 路确定性可行启发式，无值函数、无前瞻。(ii) \textbf{预算不在准入处执法}：调度器源码零预算引用；ACB 的 token/成本/墙钟预算是策略元数据，未被任何路径强制；唯一强制的预算是系统调用授权处的两阶段 reserve/commit（$\chi_9$）。若要忠实建模，约束必须迁到授权边界或扩展准入状态。(iii) \tcode{costUsdMicros} 通道端到端存在但 LLM 门面只报告 token——成本约束在实践中不可消耗。这三条是\emph{现状的}工程不变式；CMDP 是把“应该是什么”形式化的\emph{提议}（第~\ref{ch:future} 章研究问题 RQ2）。

% 第 10 章：Repository Monitoring 与 Event-Triggered Execution
\chapter{仓库监控与事件触发执行}\label{ch:monitoring}

\section{监督子系统的真实形状（轮询，而非 Webhook）}

仓库监督是\textbf{轮询模型}（源码核验，\tcode{heartbeat/repositorySupervisor.ts}）：每仓库 \tcode{setInterval(pollIntervalMs)} 驱动一次 \tcode{scanRuntime}$\to$\tcode{probe.observe()}（本地 Git 探针）；重叠扫描被生命周期令牌丢弃。每次观测计算\textbf{观测摘要}
\begin{equation}
\mathrm{digest}_t=\mathrm{SHA256}\bigl([\,headSha \,\|\, \text{UNBORN},\;\text{workflowDigest},\;\text{normalizedChangedFiles}\,]\bigr),
\end{equation}
摘要不变 $\Rightarrow$ 无待处理变更；变化则武装一个 \tcode{debounceMs} 计时器，静默后冲刷（flush）一个\emph{去重}的 \tcode{RepositoryEvent}，其身份 dedupeKey 为常量前缀加 $\mathrm{SHA256}[\text{repoId},\,head,\,\text{workflowDigest}]$。工作流摘要覆盖启用的自动化与 \tcode{on\_change} 绑定：绑定集变化会重新武装事件（摘要变 $\Rightarrow$ 新键），而未启用绑定的仓库保持字节相同摘要（无“事件风暴”的兼容性设计）。失败诚实降级（\tcode{degraded} 状态 + 降级脉冲）。监督\textbf{只观测与记录}，不自行入队；审计执行桥同样经规范化路径（\tcode{launchDefinitionRun}）以围栏式 \tcode{created}$\to$\tcode{queued} 状态守卫单次尝试排空。

\section{事件触发 vs 周期触发：结构性类比（\STRONG）}

该机制的形态是“Riemann 采样观察下的 Lebesgue 事件”：轮询 + 摘要比较是\emph{周期观测者}；摘要变化是\emph{事件}；防抖是\emph{驻留时间}（dwell time / 最小事件间隔）。事件触发采样优于周期采样的经典结论~\cite{astrom1999} 在结构上支持“变更触发评审在同等最坏时延与动作预算下浪费更少运行”这一\emph{定性}主张；Tabuada 的最小事件间隔时间（minimum inter-event time）形式化了防抖作为防抖振（Zeno 行为）手段的角色~\cite{tabuada2007}。\textbf{但}两文的定量定理都以随机被控对象与噪声模型为前提——Git 仓库在提交之间\emph{没有动力学}：它是一个被无噪声观测的分段常数信号。因此：\textbf{结构映射 \STRONG；最优性内容 \WEAK}。任何“变更触发严格优于 cron 触发 $x\%$”的定量主张在本报告的框架内都是未证明的（第~\ref{ch:future} 章研究问题 RQ5）。

\section{CUSUM/变化检测：不适用（负结果，\NA）}

顺序变化检测理论（Page 的连续检查方案~\cite{page1954} 及其后续）存在的理由是\emph{噪声观测}下的检测延迟/虚警权衡。摘要比较是无噪声、内容寻址的哈希等值测试：虚警概率 $\approx$ 碰撞概率 $\approx 0$，漏检概率 $=0$，因此\emph{没有阈值可优化、没有权衡可言}——它只在与 \tcode{a!==b} 相同的意义上是“变化检测”。把防抖称作 CUSUM 阈值正是本研究要防止的“数学优雅强迫映射”。变化检测理论在\emph{观测变为有噪声}时才变得相关：远程镜像的同步滞后、部分写入、基于 mtime 的启发式、CI 指标的漂移检测——这些是假设性的未来通道，记录为条件性映射而非现状描述。

\section{单运行去重与“缓冲大小为 1”}

监督的去重 + 触发计划的“每 (repository, definition, event) 至多一次运行”共同实现了\emph{每键缓冲大小为 1、新到替换旧}的策略：突发的连续提交不会排队成一串过时的评审运行，而是被合并为最新状态的至多一个待执行计划。信息年龄文献的奠基性洞察——对新鲜度而言“小缓冲 + 丢弃过时更新”优于“无界排队”~\cite{kaul2012}——正是该设计的理论理由（这是 AoI 理论向本系统\emph{真正}迁移的一条；其余部分的诚实评估见第~\ref{ch:freshness} 章）。

% 第 11 章：Evidence、Provenance 与 Reproducibility
\chapter{证据、溯源与可复现性}\label{ch:evidence}

\section{证据模型}

证据记录为强类型 11 元组：
\begin{align}
e&=\bigl(\text{source},\,\text{ruleId},\,\text{loc}(\text{path},\text{startLine},\text{endLine}),\nonumber\\
&\qquad\text{conf}\in[0,1],\,\text{payload},\,\mathrm{prov}\bigr),\\
\mathrm{prov}&=(\text{repository},\text{sha},\text{analyzer},\text{analyzerVersion}),
\end{align}
$\text{source}\in\{\text{ast},\text{sast},\text{git},\text{lint},\text{symbol},\text{test},\text{agent}\}$。身份函数为
\begin{equation}
\fp(e)=\mathrm{SHA\text{-}256}\bigl(\canon(N(e))\bigr)\in\{0,1\}^{256},
\end{equation}
其中 $N$ 是 store 侧的确定性规范化（路径规范化——与仓库系统调用同一 \tcode{normaliseResourcePath}，拒绝绝对路径/穿越/NUL；\textbf{store 自行计算指纹，从不信任分析器提交的指纹}；校验置信度/行号/溯源；深冻结），$\canon$ 是 11 元组的规范 JSON 序列化（递归键排序、\emph{数组保序}、深度上限 64、拒绝非纯对象与环）。两个设计点值得强调：\emph{溯源参与哈希}（同一规则在不同修订版本永不共享身份）；\emph{身份与标签分离}（$\tcode{evid}\_\langle uuid\rangle$ 是非身份标签）。

\section{理论映射}

\begin{itemize}[itemsep=0.05em]
  \item \textbf{why/where 溯源}：Buneman--Khanna--Tan 的经典区分——“结果由哪些源数据贡献（why）”与“值从哪个源位置拷贝（where）”~\cite{buneman2001}——与“发现必须引用贡献证据（why）+ 证据钉定 (path, 行区间)（where）”一一对应。映射强度：\STRONG（形式场景是只读关系查询，非仓库分析，但两个溯源问题逐字迁移）。
  \item \textbf{PROV 词汇}：$\mathrm{prov}$ 四元组是 W3C PROV-DM 实体/活动/代理/派生词汇的实例化——证据（实体）由分析器运行（活动）使用快照（实体）生成~\cite{provdm2013}。\DIRECT（结构层）。
  \item \textbf{内容寻址承诺}：SHA-256-规范元组指纹是 Merkle 哈希树的单叶退化情形~\cite{merkle1989}；RepositorySnapshot 的对象库读取同属内容寻址谱系。\DIRECT（机制层）。
  \item \textbf{审计日志}：仅追加、全决策入账、永不出现原始句柄的 AuditJournal 站在安全审计日志~\cite{schneier1999} 与防篡改日志数据结构（history tree 及其成员/一致性证明）~\cite{crosby2009} 的谱系上。\STRONG（本系统对手模型是内部一致性而非事后取证妥协，声明更弱；一致性证明列为未来升级路径）。
  \item \textbf{供应链证明}：分析器（步骤）绑定输入（钉定快照=材料）到输出（证据=产品）、步骤身份=分析器+版本的形态与 in-toto 的 layout+link 元数据同构~\cite{torresarias2019}。\STRONG。
  \item \textbf{模型输出的事后接地}：评审 Agent 的输出在被确定性接地检查（接受/降级/拒绝）之前是\emph{不可信}的——这是 RARR 式“检索证据并修订输出直到被支撑”~\cite{gao2023rarr} 的内核化版本：接地依据不是另一个模型，而是第~\ref{ch:security} 章治理下的确定性证据库。实证动机：原始模型审查输出的行动率类型依赖且不可靠~\cite{goldman2025}。
\end{itemize}

\section{接地不变式（两级严格度，如实）}

\begin{proposition}[证据接地]\label{prop:grounding}
(a) 工作流运行时层面：schema 强制 $\forall f\;\exists E_f\subseteq E_t,\;E_f\neq\varnothing$（\tcode{evidenceIds} 数组 \tcode{min(1)}）。(b) 评审工作负载的接地函数 $\ground:\mathrm{Finding}^{*}\to\{\text{accept},\text{downgrade},\text{reject}\}^{*}$ 是确定性分类器：模型引用未知证据 id $\Rightarrow$ 拒绝；未变更文件引用 $\Rightarrow$ 拒绝；越界行号 $\Rightarrow$ 拒绝；\tcode{confirmed} 主张落在变更 hunk 之外或无确定性信号 $\Rightarrow$ 降级为 \tcode{likely}；被接受的发现被确定性附加运行证据。(c) \textbf{两级严格度}：遗留评审 schema 层（迁移期）\tcode{evidenceIds} 可选——finding \emph{可以}在 schema 层无证据存在；仅工作流运行时层强制 $\ge1$。
\end{proposition}
\tagm{(a)(b) ENGINEERING\_INVARIANT；(c) 记录为限制}

\begin{figure}[htbp]
\centering
\begin{adjustbox}{max width=\textwidth}
\begin{tikzpicture}[
  nd/.style={draw=framegray,fill=boxgray,rounded corners=2pt,font=\scriptsize,align=center,minimum height=2.0em,inner sep=4pt},
  sn/.style={draw=accentblue!70,fill=softblue,rounded corners=2pt,font=\scriptsize,align=center,minimum height=2.0em,inner sep=4pt},
  ev/.style={draw=framegray,fill=softamber,rounded corners=2pt,font=\scriptsize,align=center,minimum height=2.0em,inner sep=4pt},
  fd/.style={draw=accentred!70,fill=softrose,rounded corners=2pt,font=\scriptsize,align=center,minimum height=2.0em,inner sep=4pt},
  ln/.style={-{Stealth[length=2mm]},inkgray!80},
  lb/.style={font=\tiny\itshape\color{inkgray},midway,fill=white,inner sep=1pt}]
\node[sn] (git) at (0,2.4) {Git 对象库\\$headSha$};
\node[nd] (snap) at (3.4,2.4) {RepositorySnapshot\\（不可变钉定）};
\node[nd] (an1) at (6.8,3.4) {分析器 $A_{v_1}$\\style/secret v1.0.0};
\node[nd] (an2) at (6.8,1.4) {分析器 $A_{v_2}$\\（另一版本）};
\node[ev] (e1) at (10.2,3.4) {证据 $e_1$：$\fp=\mathrm{SHA256}(\canon)$};
\node[ev] (e2) at (10.2,1.4) {证据 $e_2$};
\node[fd] (f1) at (13.4,2.4) {发现 $f$（confirmed）\\$E_f\ni e_1$};
\node[nd] (aud) at (3.4,0.0) {AuditJournal\\（仅追加）};
\draw[ln] (git) -- (snap) node[lb,above]{rev-parse 验证};
\draw[ln] (snap) -- (an1) node[lb,above]{读取};
\draw[ln] (snap) -- (an2) node[lb,below]{读取};
\draw[ln] (an1) -- (e1) node[lb,above]{$\mathrm{prov}.sha, v_1$};
\draw[ln] (an2) -- (e2) node[lb,below]{$\mathrm{prov}.sha, v_2$};
\draw[ln] (e1) -- (f1) node[lb,above]{接地 $\ground$};
\draw[ln,dashed,inkgray!60] (e2.south) to[bend right=14] node[lb,below]{无确定性信号 $\Rightarrow$ 降级/拒绝} (f1.south);
\draw[ln,dashed,inkgray!60] (aud.north) -- (snap.south) node[lb,right]{全决策入账};
\end{tikzpicture}
\end{adjustbox}
\caption{证据溯源图。快照（蓝）由对象库钉定；分析器带版本读入；证据身份含 $(sha, analyzer, version)$；发现经确定性接地分类器 $\ground$ 引用证据（红）。虚线标注被 $\ground$ 排除的路径。}
\label{fig:provenance}
\end{figure}

\section{确定性与指纹可复现性}

\begin{assumptioninner}[{分析器确定性（版本钉定）}]\label{ass:a6}
$A_v$ 是钉定版本 $v$ 上的全函数、确定性（字节稳定输出；全性经孤立代理对替换为 U+FFFD 维持；输出按 (path, startLine, endLine, ruleId) 确定排序）——代码可证。
\end{assumptioninner}
\begin{assumptioninner}[{规范序列化的确定性}]\label{ass:a7}
JSON 数字的渲染平台稳定（IEEE-754 双精度在 TS 引擎间的确定性渲染）——\emph{工程}假设。
\end{assumptioninner}
\begin{assumptioninner}[{规范化的单射性}]\label{ass:a7p}
$\canon\circ N$ 在值域上单射（不同证据内容产生不同规范串）。\textbf{已知边界}：深度上限 64 的截断与数字渲染合并是\emph{确定性}碰撞通道，位于随机预言机模型之外——本假设把它们声明为工程上可忽略，而非数学上不存在。
\end{assumptioninner}
\begin{assumptioninner}[{版本可判定}]\label{ass:a8}
“相同 $v$”可从记录本身判定（$analyzerVersion$ 在 $\mathrm{prov}$ 内）。
\end{assumptioninner}

\begin{proposition}[确定性]\label{prop:determinism}
在 A6--A8 下，$(x,v)=(x',v')\Rightarrow \fp(N(A_v(x)))=\fp(N(A_v(x')))$。
\end{proposition}
\begin{proof}
$A_v$ 确定性 $\Rightarrow$ 结构化输出相等；$N,\canon,\mathrm{SHA256}$ 是函数 $\Rightarrow$ 哈希相等（函数复合保持相等；不涉及概率）。
\end{proof}
\tagm{PROPOSITION（A6 代码可证；A7/A7$'$ 工程假设）}

\begin{proposition}[碰撞界（随机预言机理想化）]\label{prop:collision}
把 SHA-256 建模为\textbf{随机预言机}（显式声明：SHA-256 是固定公开函数，RO 是仅为得出界而用的建模虚构）。设 $n$ 条证据记录的\textbf{规范编码两两不同}（相同元组的重复是去重而非碰撞），$b=256$，则
\begin{equation}
\Pr\bigl[\exists\, i<j:\ \fp_i=\fp_j\bigr]\;\le\;\binom{n}{2}2^{-b}\;\le\;\frac{n^{2}}{2^{\,257}}.
\end{equation}
\end{proposition}
\begin{proof}
逐对使用 union bound；RO 模型下每个\emph{不同}输入对的碰撞概率为 $2^{-256}$。注意 union bound \emph{不需要}碰撞事件间的独立性——只需要逐对界。数值：$n=10^{9}\approx2^{30}$ 时上界 $\le 2^{-197}$。
\end{proof}
\tagm{THEOREM\_FROM\_KNOWN\_MODEL（RO 理想化下的 union/birthday 界；Merkle 谱系~\cite{merkle1989}）}

\section{真相不变式 T1/T2 与可复现性 R1}

\begin{proposition}[真相不变式]\label{prop:truth}
\textbf{T1（Git 事实来自 Git 对象库）}：快照读取是纯函数 $\mathrm{Read}:(repo,sha,path)\to blob$，经失败关闭的 \tcode{rev-parse --verify} 钉定；后续工作树变更不能改变既有快照的观测（不可变性来自对象库寻址本身）。\textbf{T2（PR 事实来自提供商 API）}：外部 PR 状态只经提供商中介、失败关闭的路径进入（未配置 $\Rightarrow$ 类型化错误）；密钥仅服务端；发布只经授权意图。
\end{proposition}
\tagm{ENGINEERING\_INVARIANT（T1 逐行代码可证；T2 部分代码接地、部分模型级补全；提供商诚实性在模型之外）}

\begin{proposition}[可复现性 R1]\label{prop:repro}
钉定快照 $+$ 确定性分析器 $\Rightarrow$ 相等输入产生相等证据指纹（命题~\ref{prop:determinism}）；碰撞概率可忽略（命题~\ref{prop:collision}，RO 理想化 + A7$'$）。这是可复现构建（reproducible-builds）原则在\emph{分析}侧的实例化：同一源、同一过程版本 $\Rightarrow$ 位相同输出。
\end{proposition}
\tagm{PROPOSITION（输入侧 ENGINEERING\_INVARIANT）}

\paragraph{两个风险世界的如实声明。}
当前存在\emph{两个}风险评分世界：遗留 Python 引擎的 \tcode{compose\_file\_risk}（信号分数的加权钳制：$0.28\,$style $+\,0.39\,$structural $+\,0.33\,$semantic $+\,0.05\,$dup $+\,0.5\,$security，带安全地板）与内核证据置信度——前者的权重\emph{不是}从后者推导的。任何“风险评分有理论保证”的说法都不成立；其公理化（单调性、校准）是开放问题（第~\ref{ch:future} 章研究问题）。

% 第 12 章：Repository / Evidence Freshness
\chapter{仓库与证据新鲜度}\label{ch:freshness}

本章是“怀疑式文献审计”的产物：我们专门评估了信息年龄（Age of Information, AoI）理论是否适合描述 ConsistenCy 的仓库审查新鲜度，结论是\textbf{不适用（有窄例外）}——并给出更适合的度量。该结论与正面结果同等重要。

\section{为什么线性时间年龄是错误的原始度量（负结果）}

AoI 定义时间年龄 $\Delta(t)=t-u(t)$（$u(t)$ 为最近更新时刻）~\cite{yates2021}。\textbf{休眠反例}：仓库 90 天无提交且分析钉定在 HEAD——$\Delta=90$ 天，而版本滞后 $D=0$，每条发现都完美新鲜。反过来：分析后一小时内涌入 50 个琐碎提交——$\Delta\le1$ 小时而 $D=50$，发现可能引用已不存在的代码。在低变更率区间，时间年龄与真实新鲜度\emph{负相关}；在高变更率区间它没有信息量。结论：\textbf{拒绝} $\Delta(t)$ 作为新鲜度度量；它只保留“调度时延描述符”的身份。

\section{版本滞后（\PROPOSED 度量，VAoI 接地）}

\begin{definition}[版本滞后]\label{def:versionstaleness}
$D(t)=d\bigl(\mathrm{analyzedSha}(t),\,\mathrm{HEAD}(R_t)\bigr)$，$d(s_1,s_2)=\bigl|\mathrm{anc}(s_2)\setminus\mathrm{anc}(s_1)\bigr|$（从 HEAD 可达而不从被分析提交可达的提交数）。
\end{definition}
\textbf{两个精确警示}：(i) $d$ 是\emph{非对称拟度量}——不满足三角不等式，含义是“HEAD 有而分析没有的提交”；$D$ 小\emph{不}排除“孤儿分支分析”（分析包含 HEAD 已不包含的提交——“引用不存在的代码”失败模式只被部分捕捉）；分叉历史上 $d$ 的解释需要合并基约定（线性历史上即提交计数）。(ii) 监督的“版本”实际是\textbf{摘要等价类} $[headSha, workflowDigest, changedFiles]$——重配置同 sha 也是新版本；$D$ 是它在线性历史上的影子。

该度量的形状并非新造：\emph{版本}信息年龄（Version AoI）——监控者落后源的\emph{版本计数}——正由 Buyukates 等人的工作线建立~\cite{buyukates2021}。映射判定：定义层 \USEFUL（已有文献）；把 VAoI 实例化到“评审证据监督于 Git DAG 之上”= \PROPOSED 扩展（Gossip 网络分析机件不适用）。

\section{检测时延与新鲜证据时刻（初等界）}

\begin{assumptioninner}[{新鲜度界的确定性前提}]\label{ass:a9}
变更时刻外生；扫描不重叠（$\delta_{\mathrm{scan}}\le T_{\mathrm{poll}}$；重叠扫描被丢弃）；触发变更之后\emph{不再发生}摘要变化（无再武装链）。
\end{assumptioninner}

\begin{proposition}[新鲜证据时刻的上界]\label{prop:freshbound}
A9 下，时刻 $\theta$ 的变更最迟在 $\theta+T_{\mathrm{poll}}+\delta_{\mathrm{scan}}$ 被\emph{检测}（下一个轮询栅格点 + 扫描时长；摘要比较无噪声——这正是第~\ref{ch:monitoring} 章 CUSUM 不适用判定的事实基础~\cite{page1954}），再经 $T_{\mathrm{deb}}$ 后被\emph{冲刷}（防抖恰在最后一个变更后 $T_{\mathrm{deb}}$ 冲刷——A9 使其精确），再经 $T_{\mathrm{run}}$ 后有\emph{终态运行}（命题~\ref{prop:execsound} 的 F1--F3）。故对固定时刻 $t$，尾随窗口 $W:=T_{\mathrm{poll}}+\delta_{\mathrm{scan}}+T_{\mathrm{deb}}+T_{\mathrm{run}}$ 之外的滞后不可能存在：
\begin{equation}
D(t)\;\le\;N\bigl((t-W,\,t]\bigr)\quad(\text{提交计数过程})。
\end{equation}
\end{proposition}
\begin{proof}
逐段复合三个时延不等式 + 命题~\ref{prop:execsound}；窗口包含关系来自“变更至多等待 $W$ 即被终态运行覆盖”。
\end{proof}
\tagm{PROPOSITION（A9 下精确；A9 违背——再武装链、扫描重叠、崩溃（单次尝试 $\Rightarrow$ 时延无界）——即其失败条件）}

\begin{figure}[htbp]
\centering
\begin{adjustbox}{max width=0.99\textwidth}
\begin{tikzpicture}[x=1cm,y=1cm,
  ev/.style={circle,fill=accentred!80,inner sep=1.6pt},
  ln/.style={thick,inkgray}]
\draw[->,thick] (-0.2,0) -- (13.4,0) node[below right,font=\small] {$t$};
% poll grid
\foreach \x in {0.5,1.7,2.9,4.1,5.3,6.5,7.7,8.9,10.1,11.3,12.5} \draw[inkgray!35] (\x,-0.12) -- (\x,0.12);
\node[font=\tiny\itshape\color{inkgray}] at (6.5,0.42) {轮询栅格（间隔 $T_{\mathrm{poll}}$）};
% change epoch
\node[ev] (theta) at (2.3,0) {};
\node[font=\scriptsize\bfseries,above=1pt of theta] {$\theta$：变更};
% detection
\node[ev,fill=accentblue!80] (det) at (4.1,0) {};
\node[font=\scriptsize,above=1pt of det] {检测（$\le\theta+T_{\mathrm{poll}}+\delta_{\mathrm{scan}}$）};
\draw[ln,decorate,decoration={brace,amplitude=3.5pt},accentblue!70] (2.3,-0.55) -- (4.1,-0.55) node[midway,below=4pt,font=\tiny\color{inkgray}] {$\le T_{\mathrm{poll}}+\delta_{\mathrm{scan}}$};
% debounce
\node[ev,fill=accentblue!60] (fl) at (5.3,0) {};
\node[font=\scriptsize,above=1pt of fl] {冲刷};
\draw[ln,decorate,decoration={brace,amplitude=3.5pt},accentblue!50] (4.1,-0.95) -- (5.3,-0.95) node[midway,below=4pt,font=\tiny\color{inkgray}] {$T_{\mathrm{deb}}$（驻留时间）};
% run
\node[fill=softamber,draw=framegray,rounded corners=2pt,font=\scriptsize,minimum height=1.2em,inner sep=3pt] (run) at (7.3,0) {运行 $T_{\mathrm{run}}$（单次尝试）};
\node[ev,fill=softgreen!70,draw=framegray] (fin) at (9.1,0) {};
\node[font=\scriptsize,above=1pt of fin] {终态 $\Rightarrow$ 新鲜证据};
\draw[ln,decorate,decoration={brace,amplitude=3.5pt},inkgray] (2.3,-1.45) -- (9.1,-1.45) node[midway,below=4pt,font=\scriptsize\itshape] {$W=T_{\mathrm{poll}}+\delta_{\mathrm{scan}}+T_{\mathrm{deb}}+T_{\mathrm{run}}$：命题~\ref{prop:freshbound} 的尾随窗口};
% version staleness step function above
\draw[thick,accentred!80] (0.5,1.6) -- (2.3,1.6) (2.3,1.6) -- (2.3,2.05) (2.3,2.05) -- (4.1,2.05) (4.1,2.05)--(4.1,2.5) (4.1,2.5) -- (5.3,2.5) (5.3,2.5)--(5.3,2.9) (5.3,2.9) -- (9.1,2.9) (9.1,2.9) -- (9.1,1.6) (9.1,1.6) -- (12.5,1.6);
\node[font=\scriptsize\bfseries\color{accentred!80},anchor=west] at (9.3,2.4) {$D(t)$：版本滞后};
\node[font=\tiny\itshape\color{inkgray},anchor=west] at (9.3,2.05) {终态后降回 0};
\end{tikzpicture}
\end{adjustbox}
\caption{证据新鲜度时间线（\textbf{理论}示意图）。变更 $\theta$ 经检测（轮询栅格+扫描）、防抖（驻留时间）、运行后产生新鲜证据；$D(t)$ 在此期间阶梯上升、终态后降回。所有量为理论参数，非实测。}
\label{fig:freshness}
\end{figure}

\section{平均值与低速率近似（\PROPOSED 度量/推导）}

\begin{assumptioninner}[{泊松变更（随机化前提）}]\label{ass:a10}
变更时刻是速率 $\lambda_c$ 的泊松过程，且与轮询栅格独立。
\end{assumptioninner}

在 A10 下，变更的栅格相位在 $[0,T_{\mathrm{poll}})$ 上均匀（泊松增量平稳 + 与栅格独立；无检查悖论），故 $\E[\text{等到下一轮询}]=T_{\mathrm{poll}}/2$。\textbf{冲刷参照}的平均检测-冲刷时延为
\begin{equation}
\E[\text{delay}_{\mathrm{flush}}]\;=\;\tfrac{T_{\mathrm{poll}}}{2}+\delta_{\mathrm{scan}}+T_{\mathrm{deb}},
\end{equation}
（检测参照则为 $T_{\mathrm{poll}}/2+\delta_{\mathrm{scan}}$；两处都必须显式选择参照——本报告统一采用冲刷参照）。由更新报酬定理~\cite{ross1970}，低速率时长期平均等于逐变更均值。

\paragraph{固定时刻的期望滞后（A9 与 A10 的交界，如实表述）。}
由命题~\ref{prop:freshbound} 与泊松平稳增量：\textbf{A9 成立时}（确定性世界），$\E[D(t)]\le\lambda_c W$ 在固定 $t$ 处精确。\textbf{在 A10 下}，A9 的“无后续变更”以概率 $1-\mathrm{e}^{-\lambda_c T_{\mathrm{deb}}}>0$ 被违背——再武装链使尾随窗口包含关系失守，正确的低速率近似是
\begin{equation}
\E[D(t)]\;\lesssim\;\lambda_c\bigl(W+\E[\text{再武装扩散}]\bigr),\qquad \text{当 } \lambda_c\,(T_{\mathrm{poll}}+T_{\mathrm{deb}})\ll1\text{ 时误差为 }O\bigl(\lambda_c(T_{\mathrm{poll}}{+}T_{\mathrm{deb}})\bigr),
\end{equation}
再武装暴露窗口延伸到下一个扫描时刻（故条件用 $T_{\mathrm{poll}}+T_{\mathrm{deb}}$ 而非仅 $T_{\mathrm{deb}}$）。\textbf{明确不主张}“峰值滞后”的期望界：无界时域上窗口计数的运行最大值几乎必然无界，任何此类主张在 A10 下都是假的；峰值陈述只允许在声明的有限时域 $H$ 上给出。

\section{AoI 迁移的最终账目}

\begin{longtable}{@{}p{0.30\textwidth}p{0.16\textwidth}p{0.46\textwidth}@{}}
\toprule
\textbf{AoI 组成部分} & \textbf{判定} & \textbf{理由}\\
\midrule
\endhead
线性/峰值时间年龄 & \NA & 休眠反例；版本离散状态。\\
缓冲-大小-1 洞察 & \USEFUL & “单运行去重+新到替换”的真实理由~\cite{kaul2012}。\\
峰值年龄记账 & \USEFUL & 检测时延+运行时间的初等界（命题~\ref{prop:freshbound}）。\\
版本 AoI & \USEFUL/\PROPOSED & $D(t)$ 的定义接地~\cite{buyukates2021}；Git-DAG/摘要等价类两个皱褶是真正的扩展。\\
事件触发 vs 周期 & \STRONG(结构) & 驻留时间框架~\cite{tabuada2007}；无定理迁移（无随机被控对象）~\cite{astrom1999}。\\
更新或等待最优采样 & \WEAK & 观测免费且精确时退化为“总是更新”；只有给运行定价后才非平凡。\\
信道/队列分析机件 & \NA & 无物理信道、无争用；拉取观测给出确定性界。\\
CUSUM/变化检测 & \NA & 哈希等值无虚警/延迟权衡~\cite{page1954}。\\
\bottomrule
\caption{AoI/新鲜度理论向 ConsistenCy 的迁移账目（怀疑式审计结论）。}
\label{tab:aoi}
\end{longtable}

% 第 13-14 章：统一数学模型 + 形式化性质
\chapter{ConsistenCy 统一数学模型}\label{ch:unified}

\section{统一模型（\PROPOSED：透镜，非实现）}

把前述各部件装配为一个离散时间演化系统：
\begin{equation}
X_{t+1}=\Phi(X_t,u_t,\zeta_t),\qquad u_t=\pi(X_t),
\end{equation}
其中 $X_t=(R_t,S_t,Q_t,\mathfrak{A}_t,C_t,V_t,K_t,E_t,P_t,B_t)$ 为全局状态（第~\ref{ch:arch} 章），$u_t$ 为准入/调度控制，$\zeta_t$ 为外生噪声（仓库变更、提供商时延/停机、崩溃事件）。\textbf{硬约束}（任何目标不得交换）：(i) 授权不变式 S1--S3；(ii) 预算不变式 $used_c+reserved_c\le max_c$（两阶段，\emph{在授权处}执法——注意与“准入处无预算执法”的现状差距，第~\ref{ch:queueing} 章）；(iii) 依赖可用性（节点仅当 DAG 前驱终态时执行）；(iv) 证据有效性（发布发现满足接地，工作流运行时层）。目标\emph{草图}：
\begin{equation}
\min_\pi\;\E\Bigl[\textstyle\sum_t\bigl(w_{\mathrm{lat}}\,\text{latency}_t+w_{\mathrm{stale}}\,g(D(t))+w_{\mathrm{cost}}\,\text{cost}_t+w_{\mathrm{risk}}\,\text{risk}_t\bigr)\Bigr]\quad\text{s.t. S1--S3 与资源约束},
\end{equation}
其中 $g$ 为未校准的滞后代价函数、风险聚合的公理未建立（第~\ref{ch:future} 章）。

\textbf{必须显式声明}：该统一模型是一个\emph{形式透镜}。实现的 $\pi$ 是 3 路确定性可行启发式（非优化器）；没有任何组件计算 $\Phi$ 或上述目标。它的价值是把“调度/新鲜度/成本/风险之间的权衡”从口头讨论提升为有明确自由度的数学问题，并为第~\ref{ch:future} 章的研究问题提供共享记号。

\paragraph{非统一性自白（non-unification confession）。}
复查全报告的命题可以发现：\emph{没有任何一个命题同时联合使用多个状态分量}——S1/S2 只读 $(C_t,B_t)$；正交性命题只对照 $(C_t)$ 与 $(K_t,V_t)$；排队节只读 $Q_t$；新鲜度节只读 $R_t$ 与 $E_t$ 的时戳；证据命题只读 $E_t,S_t$。因此本报告所称的“统一模型”应更准确地读作\textbf{共享记号（shared notation）}：一个让各章局部命题可以无歧义互相引用的公共状态空间，而不是一个产生了跨分量联合定理的模型。真正的联合优化（例如同时约束 $Q_t$ 与 $D(t)$ 的策略类）是第~\ref{ch:future} 章 RQ1/RQ2 的\emph{待做}内容，不是本报告已完成的内容。我们把这一点写在明处，以对抗“把章节摘要符号化后称作统一模型”的过度包装。

\begin{figure}[htbp]
\centering
\begin{adjustbox}{max width=\textwidth}
\begin{tikzpicture}[
  bx/.style={draw=framegray,fill=boxgray,rounded corners=2pt,font=\scriptsize,align=center,minimum height=2.2em,inner sep=4pt},
  hz/.style={draw=accentred!75,fill=softrose,rounded corners=2pt,font=\scriptsize\bfseries,align=center,minimum height=2.2em,inner sep=4pt},
  ob/.style={draw=accentblue!70,fill=softblue,rounded corners=2pt,font=\scriptsize\itshape,align=center,minimum height=2.2em,inner sep=4pt},
  ln/.style={-{Stealth[length=2mm]},inkgray!80,thick},
  lb/.style={font=\tiny\itshape\color{inkgray},midway,fill=white,inner sep=1pt}]
\node[bx,minimum width=3.1cm] (X) at (0,0) {状态 $X_t$\\$R,S,Q,\mathfrak{A},C,V,K,E,P,B$};
\node[bx] (Phi) at (4.1,0) {演化 $\Phi(\cdot,u_t,\zeta_t)$};
\node[bx] (pi) at (4.1,-1.9) {策略 $\pi$（现状：3 路启发式）};
\node[hz,minimum width=3.4cm] (cons) at (8.6,-1.9) {硬约束\\S1--S3 $\cdot$ 预算 $\cdot$ 依赖 $\cdot$ 接地};
\node[ob,minimum width=3.3cm] (obj) at (8.6,0) {目标（草图）\\时延 $+g(D(t))$ + 成本 + 风险};
\node[bx] (zeta) at (4.1,1.9) {外生 $\zeta_t$\\变更/停机/崩溃};
\draw[ln] (X) -- node[lb,above]{} (Phi);
\draw[ln] (Phi) to[bend left=16] node[lb,above]{$X_{t+1}$} (X);
\draw[ln] (pi) -- node[lb,right]{$u_t$} (Phi);
\draw[ln] (zeta) -- (Phi);
\draw[ln,dashed,accentred!75] (cons) -- node[lb,below]{不可交换} (pi);
\draw[ln,dashed,inkgray!60] (obj) -- node[lb,right]{（\PROPOSED，未实现）} (pi);
\end{tikzpicture}
\end{adjustbox}
\caption{统一形式模型（\PROPOSED\_CONSISTENCY\_FORMAL\_MODEL）。红色为任何策略都不得违背的硬约束；蓝色目标仅是研究草图——实现的 $\pi$ 不优化它。}
\label{fig:unified}
\end{figure}

\chapter{形式化安全/活性性质清单}\label{ch:properties}

下表汇总全报告的性质，含假设与“已执法 vs 期望性”的区分：

\begin{longtable}{@{}>{\raggedright\arraybackslash}p{0.9cm}>{\raggedright\arraybackslash}p{3.4cm}>{\raggedright\arraybackslash}p{3.9cm}>{\raggedright\arraybackslash}p{3.3cm}>{\raggedright\arraybackslash}p{2.4cm}@{}}
\toprule
\textbf{ID} & \textbf{性质} & \textbf{形式内容} & \textbf{关键假设} & \textbf{状态}\\
\midrule
\endhead
S1 & 无授权则无受保护效果 & $\Box(\Exec\Rightarrow\Auth(m(q)){=}1)$ & A1--A3；边界：网关穿透依赖 broker 拒绝（完全中介者是 broker） & 工程不变式（调用路径代码可证；旁路不存在性是假设）\\
S2 & 撤销支配生命周期传播 & 中介点 $\succeq t_r$ 裁决恒 0，与 $K_t$ 无关；intent carve-out & A1--A5 & 命题（A4/A5 代码可证）\\
S3 & 数据面不能直接改变裁决 & 裁决面独立性（单步、固定状态，命题~\ref{prop:verdictplane}） & 数据 $\to$ 权威的唯一通道是内核侧 Ring 校验发放 & 命题\\
T1 & Git 事实来自对象库 & $\mathrm{Read}:(repo,sha,path)\to blob$ 纯函数、失败关闭钉定 & 对象库完整性（模型外） & 工程不变式\\
T2 & PR 事实来自提供商 API & 提供商中介、失败关闭、密钥服务端、发布经意图 & 提供商诚实性（模型外） & 工程不变式（部分代码接地）\\
E1 & 发现引用持久证据 & 工作流运行时层 $\forall f\,\exists E_f\neq\varnothing$；确定性 $\ground$ & schema \tcode{min(1)}；\textbf{差距}：遗留层可选 & 工程不变式（仅工作流层）\\
R1 & 可复现证据指纹 & 命题~\ref{prop:determinism}+\ref{prop:collision} & A6--A8 + A7$'$；RO 理想化 & 命题（输入侧工程不变式）\\
L1 & 准入的合格 Agent 最终推进 & 在公平派发、依赖可用、预算未竭、提供商可用、无崩溃之下 & F4/F5 为期望性；HOL 饥饿、PENDING 吸收、$\chi_9$ 拒绝、单次尝试均是其反例条件 & 命题（假设下）；截止期执法是唯一\emph{已执法}的活性邻接机制（诚实失败而非进展）\\
\bottomrule
\caption{形式化性质总表。S = Safety，T = Truthfulness，E = Evidence，R = Reproducibility，L = Liveness。}
\label{tab:properties}
\end{longtable}

\paragraph{L1 的反例条件值得展开。}
“最终推进”在下列任一情形下失败，且每一种都\emph{设计上}被接受为诚实失败而非缺陷修复目标：(i) 高优先级持续饱和使低类饥饿（定理~\ref{thm:hol}）——由截止期取消兜底为“取消”这一终态；(ii) 共效从未提供使 Fiber 永久 PENDING；(iii) $\chi_9$ 预算拒绝是\emph{条件活性}的边界（安全优先于活性的直接体现）；(iv) LLM 无重试——一次修复尝试后类型化失败；(v) 崩溃后启动恢复把执行中标记为失败。换言之，ConsistenCy 的活性哲学是“\textbf{公平假设下的进展 + 假设失守时的诚实终态}”，而非尽力掩盖的可用性最大化——与第~\ref{ch:sep} 章指出的 OTP 分歧一致。

% 第 15-18 章：架构决策的理论解释 / 相关系统比较 / 理论支撑现状
\chapter{架构决策的理论解释}\label{ch:decisions}

下表把九项代表性架构决策映射到其工程理由、理论解释与数学性质，并标注证据强度。这是“工程直觉 $\to$ 可检查对象”转化的核心交付物之一。

\begin{longtable}{@{}>{\raggedright\arraybackslash}p{2.6cm}>{\raggedright\arraybackslash}p{3.0cm}>{\raggedright\arraybackslash}p{3.2cm}>{\raggedright\arraybackslash}p{3.2cm}>{\raggedright\arraybackslash}p{1.7cm}@{}}
\toprule
\textbf{架构决策} & \textbf{工程理由} & \textbf{理论} & \textbf{数学性质} & \textbf{强度}\\
\midrule
\endhead
逐 syscall 授权（9 步链重评） & 防缓存授权过期、防 TOCTOU、撤销即时生效 & 引用监视器/完全中介~\cite{anderson1972,saltzer1975,bishop1996} & S1：$\Box(\Exec\Rightarrow\Auth)$ & \DIRECT\\
撤销永久单调 + 下一调用拒绝 & 简化撤销语义；安全不依赖传播时序 & 撤销经中介（退化情形）~\cite{redell1974,miller2003} & S2：支配性（命题~\ref{prop:revocation}） & \DIRECT\\
Coeffect 与授权两次独立检查 & 可用性与权威分属不同时间尺度 & vat 模型/事件循环响应式~\cite{miller2006} & 正交性（命题~\ref{prop:orthavail}、\ref{prop:propagation}） & \DIRECT\\
Context VM 与能力分离 & “看见什么”与“能做什么”不同轴 & 关注点分离/最小特权~\cite{saltzer1975} & $\Vis\nvdash\Auth$（命题~\ref{prop:orthvis}） & \DIRECT\\
intent 分发门（不可逆效果） & 外部副作用需要幂等与显式意图 & 幂等生产者/效果分级的业务效果~\cite{debenedetti2025} & $\mathrm{dispatch(intent)}\Rightarrow\neg$内联（\S\ref{ch:workflow}） & \STRONG\\
Scheduler 属于 Kernel（非插件） & 准入/并发/截止期是权威而非功能 & 准入控制 + 非抢占 HOL~\cite{kleinrock1975} & 确定性准入函数；饥饿形状（定理~\ref{thm:hol}） & \DIRECT{}（词汇）\\
Evidence-first findings & 模型输出不可信，需事后接地 & why/where 溯源 + 事后归因~\cite{buneman2001,gao2023rarr} & $\ground$ 分类器；E1 & \STRONG\\
RepositorySnapshot SHA 钉定 & 读取必须是内容函数 & 内容寻址承诺~\cite{merkle1989} & T1；$A_v$ 复合（R1） & \DIRECT\\
持久触发计划 + 去重键 & 恰好一次发起、防事件风暴 & AoI 缓冲-1 洞察~\cite{kaul2012} & 单飞 CAS 状态机（\S\ref{ch:workflow}） & \USEFUL\\
\bottomrule
\caption{架构决策 $\to$ 理论解释映射表。}
\label{tab:decisions}
\end{longtable}

\chapter{与相关系统/论文的比较}\label{ch:comparison}

\begin{longtable}{@{}>{\raggedright\arraybackslash}p{2.2cm}>{\raggedright\arraybackslash}p{2.6cm}>{\raggedright\arraybackslash}p{2.6cm}>{\raggedright\arraybackslash}p{2.5cm}>{\raggedright\arraybackslash}p{2.5cm}@{}}
\toprule
\textbf{系统} & \textbf{权威中介} & \textbf{证据/接地} & \textbf{验证程度} & \textbf{失败哲学}\\
\midrule
\endhead
\textbf{ConsistenCy v3} & 逐 syscall 9 步能力链；intent 门 & 确定性分析器 + 内容指纹 + 强制 $\ground$ & 无形式证明；模型级命题 + 源码核验 & fail-closed、单次尝试、诚实终态\\
seL4~\cite{klein2009} & 能力内核，同步中介 & N/A（机制不做证据） & \textbf{完整机器检查函数正确性证明} & 正确性由证明承担\\
Capsicum~\cite{watson2010} & 能力模式（沙箱内禁全局命名空间） & N/A & 无证明；部署经验 & 细粒度授权实践\\
in-toto~\cite{torresarias2019} & 签名 layout 授权步骤 & 链接元数据绑定材料$\to$产品 & 无证明；部署（供应链） & 步骤可验证\\
CT/history tree~\cite{crosby2009} & N/A & 防篡改仅追加日志 + 一致性证明 & 结构证明 & 公开可审计\\
Erlang/OTP~\cite{armstrong2003} & 无能力模型（进程隔离） & N/A & 无证明；工业可靠性记录 & \textbf{让它崩溃 + 重启掩盖}\\
ReAct~\cite{yao2023react} & \textbf{无}（宿主约定） & 无 & N/A & 提示层约束\\
AutoGen~\cite{wu2024autogen} & \textbf{无}（会话编程） & 无 & N/A & 应用层\\
MetaGPT~\cite{hong2024metagpt} & \textbf{无}（SOP 编码） & 无 & N/A & 应用层\\
CodeReviewer~\cite{li2022codereviewer} & N/A & \textbf{无证据接地}（模型端到端） & 实证（预训练） & 概率输出\\
\bottomrule
\caption{与相关系统/论文的比较。LLM Agent 框架一行显示的“无权威中介”正是 ConsistenCy 的存在理由；与 OTP 的失败哲学分歧（诚实失败 vs 重启掩盖）见第~\ref{ch:sep} 章。}
\label{tab:comparison}
\end{longtable}

定位声明：ConsistenCy 不是“又一个 Agent 框架”，而是把 OS 内核六十年的权限理论~\cite{dennis1966,anderson1972,saltzer1975}与可复现计算的溯源理论~\cite{buneman2001,provdm2013,merkle1989}装配到 LLM Agent 工作负载上的\emph{基础设施}；ReAct/AutoGen/MetaGPT 属于其上可承载的工作负载形态，而非同类竞争者。

\chapter{哪些理论已经支撑当前实现}\label{ch:supported}

图~\ref{fig:matrix} 以矩阵形式总结全部理论映射的强度分布：\DIRECT{} 集中在安全与分离公理（引用监视器、撤销、正交性、内容寻址），\STRONG{} 集中在生命周期/溯源/意图门，\USEFUL{} 在排队词汇与缓冲洞察，\PROPOSED{} 在 CMDP 与版本滞后度量，\NA{} 记录了 Petri 机件、CUSUM 与 AoI 线性年龄三组负结果。

\begin{figure}[htbp]
\centering
\small
\begin{tabular}{@{}l|c|c|c|c|c|c@{}}
\toprule
\textbf{概念 $\backslash$ 强度} & \DIRECT & \STRONG & \USEFUL & \PROPOSED & \WEAK & \NA\\
\midrule
逐调用授权/引用监视器 & \cellcolor{softblue}$\bullet$ & & & & & \\
撤销支配 & \cellcolor{softblue}$\bullet$ & & & & & \\
可用性/授权正交 & \cellcolor{softblue}$\bullet$ & & & & & \\
可见性/权限正交 & \cellcolor{softblue}$\bullet$ & & & & & \\
内容寻址身份 & \cellcolor{softblue}$\bullet$ & & & & & \\
ACB/Fiber LTS & \cellcolor{softblue}$\bullet$ & & & & & \\
PROV 溯源结构 & \cellcolor{softblue}$\bullet$ & & & & & \\
\midrule
Fiber 监督谱系 & & \cellcolor{softgreen}$\bullet$ & & & & \\
证据 why/where & & \cellcolor{softgreen}$\bullet$ & & & & \\
意图门/幂等 & & \cellcolor{softgreen}$\bullet$ & & & & \\
防篡改日志谱系 & & \cellcolor{softgreen}$\bullet$ & & & & \\
事件触发结构 & & \cellcolor{softgreen}$\bullet$ & & & & \\
LLM 服务时间实证 & & \cellcolor{softgreen}$\bullet$ & & & & \\
\midrule
排队记号/稳定性/Little & & & \cellcolor{softamber}$\bullet$ & & & \\
AoI 缓冲-1 洞察 & & & \cellcolor{softamber}$\bullet$ & & & \\
峰值年龄记账 & & & \cellcolor{softamber}$\bullet$ & & & \\
模型检验扩展 & & & & \cellcolor{softrose}$\bullet$ & & \\
CMDP 准入研究模型 & & & & \cellcolor{softrose}$\bullet$ & & \\
版本滞后 $D(t)$ & & & & \cellcolor{softrose}$\bullet$ & & \\
\midrule
两阶段预算预留 & & & & & \cellcolor{boxgray}$\bullet$ & \\
更新或等待最优采样 & & & & & \cellcolor{boxgray}$\bullet$ & \\
Petri/WF-net 机件 & & & & & & \cellcolor{softrose}$\bullet$\\
CUSUM/变化检测 & & & & & & \cellcolor{softrose}$\bullet$\\
线性时间年龄 & & & & & & \cellcolor{softrose}$\bullet$\\
\bottomrule
\end{tabular}
\caption{理论支持强度矩阵（图例：蓝=\DIRECT，绿=\STRONG，琥珀=\USEFUL，玫瑰=\PROPOSED/\NA，灰=\WEAK）。逐项论证与文献见第~\ref{ch:security}--\ref{ch:freshness} 章及附录~\ref{app:mapping}。}
\label{fig:matrix}
\end{figure}

\chapter{哪些只是提议研究方向}\label{ch:proposed}

以下内容\textbf{不是}当前系统的属性，把它们写成现状即构成理论包装。（评估深度声明：本章列举的负结果与新方向中，AoI/新鲜度经受了最深的专项审计（第~\ref{ch:freshness} 章），Petri 与 CUSUM 为边界论证（\S\ref{ch:workflow}、\S\ref{ch:monitoring}），深度不一；此处列出负结果的目的是划定“已确立理论”的边界，而非宣称对三者的评估等量齐观。）
\begin{enumerate}[itemsep=0.05em]
  \item \textbf{CMDP 准入优化}：现状是 3 路可行启发式；预算不在准入处执法；成本通道不可消耗。
  \item \textbf{版本滞后度量 $D(t)$ 与触发策略优化}：度量是本报告的提议（VAoI 接地）；“事件触发 vs 周期的定量优劣”是未证明命题。
  \item \textbf{积 LTS 的模型检验}：48 状态空间是现成的检验对象，但未接线。
  \item \textbf{统一模型 $\Phi$ 与目标函数}：透镜，无组件计算它。
  \item \textbf{风险评分公理化}：两世界并存的现状下，任何单调性/校准主张都未建立。
  \item \textbf{审计日志一致性证明}：仅追加接口 + 指纹纪律是现状；Crosby 式密码学防篡改~\cite{crosby2009} 是升级路径。
  \item \textbf{ensemble 可靠性模型}：没有数据支撑独立假设，禁止 $\prod_i p_i$ 式计算；贝叶斯集成仅列为经验研究议程。
\end{enumerate}

% 第 19-20 章 + 结论
\chapter{局限}\label{ch:limits}

\section{模型与方法论的局限}
\begin{enumerate}[itemsep=0.1em]
  \item \textbf{模型级而非实现级证明}。除标注“代码可证”的构造性论据外，S1--S3、正交性与活性命题证明的是架构模型 $\mathcal{M}_C$。“不存在旁路”（A1 的全称否定）只有形式化验证才能消除——seL4 正是走了这条路~\cite{klein2009}，而 ConsistenCy 没有。
  \item \textbf{OS 围堵边界}。除子进程隔离（剥离父环境秘密 + RPC 中介）外，OS 级文件系统/网络围堵\emph{未强制}；Ring 3 in-process 内建组件与宿主共享进程边界。S1 只约束经中介的效果。
  \item \textbf{外部库公理}。Fiber LTS 语义由 Cordis 4.0.0-rc.8 库拥有；库升级会重新打开第~\ref{ch:states} 章结论。
  \item \textbf{理想化}。碰撞界依赖随机预言机 + A7$'$（规范化单射为工程假设）；排队节整体是声明的基线近似（假设表~\ref{tab:qassump} 含五项违背/存疑）。
  \item \textbf{诚实差距清单}（docs$\neq$code）：授权链“8 步”注释 vs 9 步实现；预算注释声称调度器消费而源码无引用；ACB 预算元数据未执法；成本通道不可消耗；Coordinator 内存去重；网关未注册动作穿透；遗留 schema 允许无证据 finding。每条都已编入相应命题的边界条件。
  \item \textbf{文献覆盖}。45 篇经核验文献覆盖五条 lane 的正典；两篇现代平行工作（CaMeL~\cite{debenedetti2025}）为未经评审预印本，其元数据以 2026-08 为准。
\end{enumerate}

\section{度量与评估的局限}
本报告\emph{没有}任何实测基准：排队曲线是理论图（图~\ref{fig:queueing} 图注已声明）；$D(t)$、$\E[\text{delay}_{\mathrm{flush}}]$ 等度量的参数（$T_{\mathrm{poll}},T_{\mathrm{deb}},\dots$）未在真实部署上标定；风险评分两世界并存。经验评估（真实仓库、真实 LLM、真实 PR 流）是把本报告的透镜变成工程反馈的必要下一步。

\chapter{未来研究议程}\label{ch:future}

\begin{openproblem}[RQ1 高频提交下的证据滞后最小化]
在提交率 $\lambda_c$、轮询 $T_{\mathrm{poll}}$、防抖 $T_{\mathrm{deb}}$、运行成本与并发上限约束下，最小化 $\E[g(D(t))]$（或有限时域峰值）的策略类是什么？变更触发与周期触发的定量 crossover 在哪个参数区域？（\PROPOSED；需要校准的 $g$。）
\end{openproblem}
\begin{openproblem}[RQ2 token/成本约束下的多 Agent 调度]
把 $\chi_9$ 的两阶段预算提升为 CMDP 准入约束（或证明确定性启发式的竞争比）；估计 $\mu$ 的在线方法；$c>1$ 与批化到达下的重交通界。（\PROPOSED~\cite{altman1999}。）
\end{openproblem}
\begin{openproblem}[RQ3 撤销与异步生命周期的安全一致性]
陈旧门面窗口 $W_{\mathrm{stale}}$ 的有界性证明（微任务队列调度分析）；意图门 carve-out 的“恰好一次效果发起”在 Sink 持久性故障模型下的形式化。
\end{openproblem}
\begin{openproblem}[RQ4 可验证的 Agent 产出软件证据]
把 $\ground$ 从分类器扩展为论证框架（发现=主张，证据=支撑，$\ground$=攻击关系）；证据集增长下接受精度的单调性；签名/一致性证明的引入路径~\cite{crosby2009}。
\end{openproblem}
\begin{openproblem}[RQ5 事件触发评审的何时更优]
在 $g$ 校准后，更新报酬框架下变更触发 vs 周期触发的长期平均成本比较；与“驻留时间”参数 $T_{\mathrm{deb}}$ 的联合优化~\cite{astrom1999,tabuada2007}。
\end{openproblem}
\begin{openproblem}[RQ6 上下文分配作为约束信息分配]
$\Vis$ 投影在 token 预算下的最优页面集选择（工作集理论的 LLM 版）；与 $\Auth$ 的正交性保持。
\end{openproblem}
\begin{openproblem}[RQ7 Agent Harness 的长期仓库稳定性]
把“多 Agent 持续修改下仓库质量的平稳性”形式化：风险评分的公理化（单调、校准聚合）、评审漏检过程的诚实建模（无独立假设）、以及两风险世界的合并。
\end{openproblem}

\section*{结论}

本报告对 ConsistenCy v3 完成了一次完整的理论化闭环：以当前工作树为唯一事实来源，把 14 个源码核验的组件映射到六个理论域；经 45 篇网络核验文献的支持强度分级；构造了覆盖授权、撤销、正交性、生命周期、工作流、排队、证据与新鲜度的形式模型 $\mathcal{M}_C$，其中包含 2 个已确立定理的直接实例、11 个模型级命题（含全部假设、证明与失败条件）与 3 组负结果（Petri 机件、CUSUM、线性时间年龄）。敌意数学审计（PASS\_AFTER\_FIXES，全部修正已并入）与独立引文审计（零虚构，5 处修正已应用）之后，结论可以概括为三句话：
\begin{enumerate}[itemsep=0.1em]
  \item \textbf{ConsistenCy 的核心安全直觉是理论上的直系后代}：引用监视器、对象能力、撤销经中介与内容寻址承诺在机制层一一对应（\DIRECT），其公理“Coeffect $\neq$ Authorization”“Context VM $\neq$ Permission System”获得了精确的数学表述——安全只由 $\Auth$ 承担，可用性与可见性只影响活性。
  \item \textbf{最有价值的部分恰恰是诚实的负结果}：AoI 的线性年龄、Petri 网机件与 CUSUM 都被判定为不适用并给出理由；CMDP、版本滞后与模型检验被明确标记为提议而非现状。
  \item \textbf{从工程直觉到研究议程的通道已经打开}：统一模型 $\Phi$、七个开放问题与“真实测量”的缺失共同构成下一步——让透镜遇到数据。
\end{enumerate}


\appendix
% 附录 A：完整数学推导
\chapter{完整数学推导}\label{app:derivations}

本附录给出正文中压缩陈述的完整推导。所有推导针对架构模型 $\mathcal{M}_C$；随机性只在显式声明的假设（A9/A10、随机预言机）下出现。

\section{Erlang-C 公式的推导梗概（定理~\ref{thm:erlangc}）}

$M/M/c$ 的稳态是生灭链：$\pi_{n+1}=\frac{\lambda}{\min(n{+}1,c)\mu}\pi_{n}$。记 $a=A_e=\lambda/\mu$，则 $\pi_n=\pi_0\,a^n/n!\ (n\le c)$，$\pi_n=\pi_0\,a^n/(c!\,c^{\,n-c})\ (n>c)$。等待概率即 Erlang-C：
\begin{equation}
P_{\mathrm{wait}}=\sum_{n\ge c}\pi_n=\frac{a^c}{c!}\cdot\frac{1}{1-\rho}\cdot\Bigl[\sum_{k=0}^{c-1}\frac{a^k}{k!}+\frac{a^c}{c!}\cdot\frac{1}{1-\rho}\Bigr]^{-1},\qquad \rho=\frac{a}{c}<1.
\end{equation}
条件等待是速率为 $c\mu$ 的指数剩余服务 $k$ 次（$k$ 为到达时队长 $-c$），几何混合后 $\E[W_q\mid\text{wait}]=1/(c\mu-\lambda)$，故 $\E[W_q]=P_{\mathrm{wait}}/(c\mu-\lambda)$。$c=1$：$P_{\mathrm{wait}}=\rho$，$\E[W_q]=\rho/(\mu-\lambda)=\rho\cdot\tau/(1-\rho)$。到达看到稳态（PASTA）使“到达时队长分布 $=$ 稳态分布”成立——这是把 $\pi$ 读作到达分布的合法性来源。以上为 $M/M/c$ \emph{理想化}内的定理~\cite{kleinrock1975}。

\section{Little 定律的证明梗概（定理~\ref{thm:little}）}

以 $A(t)$ 记 $[0,t]$ 到达数、$D(t)$ 离去数，系统内人数 $L(t)=A(t)-D(t)$。每个顾客的累积在系统时间 $\int_0^t L(s)\,ds=\sum_{i\text{ 已离去}}W_i+\sum_{i\text{ 在系统}}(\text{部分})$。若 (i) 到达率极限 $\lambda=\lim_{t\to\infty}A(t)/t$ 存在且有限，(ii) 均值逗留 $W=\lim_{n\to\infty}\frac1n\sum_{i\le n}W_i$ 存在且有限，(iii) 顾客最终离去（无永久滞留），则由 Cauchy--Stolz 型论证 $L=\lim_{t\to\infty}\frac1t\int_0^tL=\lambda W$~\cite{little1961}。证明\emph{完全不使用}分布假设、服务纪律或优先级——这正是它在表~\ref{tab:qassump} 的违背下仍然可用的原因；代价是它只约束均值关系，不给出 $W$ 本身。

\section{非抢占 HOL 优先级均值等待（定理~\ref{thm:hol}）}

类 $k$ 顾客的等待 $=$ 残余服务 $R$（到达时刻正在服务的剩余时间）$+$ 到达时已在队的类 $\le k$ 工作量 $W_{\text{backlog}}^{(k)}$ $+$ 等待期间新到的类 $<k$ 工作量 $W_{\text{new}}^{(k)}$。均值版本（指数平滑假设下的标准分解）：
\begin{equation}
\E[W_q^{(k)}]=\E[R]+\sum_{i\le k}\lambda_i\E[W_q^{(i)}]\E[S_i]+\sum_{i<k}\lambda_i\E[W_q^{(k)}]\E[S_i],
\end{equation}
其中 $\E[R]=W_0=\sum_i\lambda_i\E[S_i^2]/2$（平均剩余服务悖论）。对 $k=1,2,\dots$ 递推求解得
\begin{equation}
\E[W_q^{(k)}]=\frac{W_0}{(1-R_{k-1})(1-R_k)},\qquad R_k=\sum_{i\le k}\rho_i
\end{equation}
（$R_k$ 为类 $\le k$ 累积利用率~\cite{kleinrock1975}。）
分母结构即饥饿形状：$R_{k-1}\uparrow1$ 时类 $k$ 均值等待 $\sim(1-R_{k-1})^{-1}$ 爆炸，即使 $R_\text{all}=\rho<1$。P3 的“严格优先级 + 类内 FIFO”是其在两类（高/低）情形的实例；ConsistenCy 的兜底是 P1 截止期取消（把无限等待转化为终态），不是公平份额~\cite{demers1989}。

\section{碰撞界的完整证明（命题~\ref{prop:collision}）}

设规范编码 $\{c_1,\dots,c_n\}$ 两两不同（A7$'$ + 前置条件），$\fp_i=H(c_i)$，$H$ 为随机预言机。对固定对 $(i,j)$：$\Pr[H(c_i)=H(c_j)]=2^{-256}$（RO 在不同输入上取值独立均匀）。则
\begin{equation}
\Pr\Bigl[\bigcup_{i<j}\{\fp_i=\fp_j\}\Bigr]\;\overset{\text{Boole}}{\le}\;\sum_{i<j}\Pr[\fp_i=\fp_j]\;=\;\binom{n}{2}\,2^{-256}\;\le\;\frac{n^2}{2}\,2^{-256}=\frac{n^2}{2^{257}}.
\end{equation}
\textbf{关键点}：Boole 不等式对事件独立性\emph{零假设}——它只需要逐对概率界，而逐对界来自 RO 在\emph{两个不同输入}上的独立取值。因此本证明没有走私任何未声明的独立性假设。数值代入 $n=2^{30}$：上界 $=2^{60}/2^{257}=2^{-197}$。

\section{确定性命题（命题~\ref{prop:determinism}）}

设 $(x,v)=(x',v')$。A6：$A_v$ 是函数且 $A_v(x)=A_v(x')$（结构化输出相等——注意这要求 A6 的\emph{字节稳定}版本而非仅仅“语义确定”）。$N$（路径规范化、校验、冻结）是函数：$N(A_v(x))=N(A_v(x'))$。$\canon$（A7 确定性渲染 + A7$'$ 单射，这里只需确定性）是函数。SHA-256 是函数。函数复合保持相等：$\fp(N(A_v(x)))=\fp(N(A_v(x')))$。$\blacksquare$。注意该证明对“同输入同输出”零概率假设；全部重量压在 A6/A7 的工程事实上，A7$'$ 只在\emph{逆命题}（不同内容 $\Rightarrow$ 不同指纹）中需要。

\section{新鲜证据时刻上界的完整证明（命题~\ref{prop:freshbound}）}

设变更（摘要变化）发生在 $\theta\in(t_g,t_{g+1}]$，$t_g$ 为轮询栅格点，栅格间隔 $T_{\mathrm{poll}}$。(1) \emph{检测}：下一次观测在 $t_{g+1}$ 开始，扫描时长 $\delta_{\mathrm{scan}}$，A9 保证不重叠（$\delta_{\mathrm{scan}}\le T_{\mathrm{poll}}$）且该次观测必看到新摘要（摘要只依赖仓库状态，观测无噪声）；故检测完成时刻 $\le t_{g+1}+\delta_{\mathrm{scan}}\le\theta+T_{\mathrm{poll}}+\delta_{\mathrm{scan}}$。(2) \emph{冲刷}：防抖计时器在检测时刻武装，A9（无后续变更）下恰经 $T_{\mathrm{deb}}$ 冲刷；冲刷时刻 $\le\theta+T_{\mathrm{poll}}+\delta_{\mathrm{scan}}+T_{\mathrm{deb}}$。(3) \emph{终态}：冲刷触发计划（去重键唯一 $\Rightarrow$ 发起一次运行），运行按命题~\ref{prop:execsound} 的 F1--F3 在有限步内到达终态；以 $T_{\mathrm{run}}$ 记其上界，新鲜证据时刻 $\le\theta+W$，$W:=T_{\mathrm{poll}}+\delta_{\mathrm{scan}}+T_{\mathrm{deb}}+T_{\mathrm{run}}$。(4) \emph{窗口包含}：任意固定 $t$，若某提交 $u\in(t-W,t]$ 且其触发链未完成，则与 (1)--(3) 矛盾；故未被终态运行覆盖的提交只能落在 $t-W$ 之前，即 $D(t)\le N((t-W,t])$。\textbf{失败条件}即 A9 的三违背：再武装链（$D(t)$ 可超窗）、扫描重叠（有效栅格变宽）、崩溃（单次尝试 $\Rightarrow$ 时延无界）。

\section{栅格相位均匀性与平均冲刷时延}

A10 下，$[0,T_{\mathrm{poll}})$ 上的相位分布：泊松过程与独立栅格的平移不变性给出 $P(\text{phase}\in ds)=ds/T_{\mathrm{poll}}$（平稳增量 + 独立性；无检查悖论——栅格不依赖变更过程）。\textbf{冲刷参照}：
\begin{equation}
\E[\text{delay}_{\mathrm{flush}}]=\E[\text{phase 的剩余等待}]+\delta_{\mathrm{scan}}+T_{\mathrm{deb}}=\frac{T_{\mathrm{poll}}}{2}+\delta_{\mathrm{scan}}+T_{\mathrm{deb}},
\end{equation}
均匀分布的期望剩余 $=T_{\mathrm{poll}}/2$。更新报酬定理~\cite{ross1970}：长期平均冲刷时延 $=$ 每周期期望总延迟/每周期期望变更数；低速率（$\lambda_c(T_{\mathrm{poll}}{+}T_{\mathrm{deb}})\ll1$，再武装概率 $1-\mathrm{e}^{-\lambda_cT_{\mathrm{deb}}}\approx\lambda_cT_{\mathrm{deb}}$ 可忽略）时逐变更均值即长期平均。\textbf{交界警示}（正文 \S\ref{ch:freshness}）：固定 $t$ 的 $\E[D(t)]\le\lambda_cW$ 只在 A9（确定性世界）内精确；A10 世界里必须加再武装扩散修正项，且无界时域“峰值”主张非法。

\section{撤销支配与传播无关性（命题~\ref{prop:revocation}、\ref{prop:propagation}）}

见正文证明；此处补全独立性论证的语义细节。$\Auth$ 的定义（定义~\ref{def:auth}）是\emph{显式给出}的公式，其自由变元集 $\mathrm{FV}(\Auth)=\{C_t,B_t,h,a,\alpha,r,s,t\}$ 由构造确定。设 $\mathcal{G}_t=(K_t,V_t,P_t,\mathfrak{A}_t,Q_t)$ 为任意其他分量組。对任意两个轨迹片段 $(X,X')$ 满足 $C=B$（逐分量）而 $\mathcal{G}\ne\mathcal{G}'$：$\Auth$ 在两轨迹上逐点相等——因为求值只读取 $\mathrm{FV}$。故“$\delta_{\mathrm{prop}}$ 任意（含 $\infty$）”与“窗口内全部拒绝”不矛盾：前者只改变 $\mathcal{G}$（Fiber/K 侧），后者由 $C$ 侧单调决定。$\blacksquare$

\section{积 LTS 良定义性（命题~\ref{prop:product}）}

$\varphi:L_{\mathrm{sched}}\rightharpoonup L_{\mathrm{fiber}}$ 部分函数且两标签集不相交。定义积迁移 $(s,f)\xrightarrow{\ell}(s',f')$ 当且仅当 (a) $\ell\in L_{\mathrm{sched}}$、$s\xrightarrow{\ell}s'$、$f=f'$，且 $\varphi(\ell)$ 未定义或与 $f$ 的共存约束一致（fiber 同步执行 $\varphi(\ell)$）；或 (b) $\ell\in L_{\mathrm{fiber}}\setminus\varphi(L_{\mathrm{sched}})$（如自主 dispose）、$f\xrightarrow{\ell}f'$、$s=s'$。每个复合标签下的迁移由分支条件互斥给出唯一性——$\varphi$ 的函数性排除“一个调度标签对应两个 fiber 标签”的歧义。不可达性：$(\tcode{NEW},\text{ACTIVE})$ 要求存在使 fiber 激活的标签共现，而 $\tcode{NEW}$ 下唯一使能标签是 register（$\varphi$ 未定义于其上），矛盾；$(\tcode{SUSPENDED},\text{ACTIVE})$ 同理由 suspend 的 $\varphi$-像 cleanup 排除。$\blacksquare$

% 附录 B：文献映射矩阵
\chapter{文献映射矩阵}\label{app:mapping}

下表为“ConsistenCy 概念 $\times$ 理论 $\times$ 关键文献 $\times$ 支持强度”的完整矩阵（正文的浓缩版是图~\ref{fig:matrix} 与表~\ref{tab:decisions}）。文献键与 \tcode{references.bib} 一致；全部 45 篇均经 2026-08-28 网络核验（研究笔记 02/06 记录了核验证据与五处修正）。

\small
\begin{longtable}{@{}>{\raggedright\arraybackslash}p{3.1cm}>{\raggedright\arraybackslash}p{2.9cm}>{\raggedright\arraybackslash}p{3.6cm}>{\raggedright\arraybackslash}p{1.5cm}>{\raggedright\arraybackslash}p{3.0cm}@{}}
\toprule
\textbf{ConsistenCy 概念} & \textbf{理论} & \textbf{关键文献} & \textbf{强度} & \textbf{备注}\\
\midrule
\endhead
默认拒绝 broker、9 步链、逐调用重评 & 引用监视器/完全中介/失败默认 & anderson1972; saltzer1975 & \DIRECT & RVM 三性质逐条对上；broker 是完全中介者（A1 边界）\\
逐调用重评（不缓存）+ sha 钉定作用域 & TOCTOU & bishop1996 & \DIRECT & 定义了所防御的漏洞类\\
不可猜测 256 位句柄；12 位指纹审计 & 对象能力不可伪造引用 & dennis1966; miller2003 & \DIRECT & 不可伪造性以熵实现（无硬件标签）\\
能力模式（默认拒绝运行时） & OS 能力模式实践 & watson2010 & \STRONG & “能力层可在通用环境落地”的先例\\
形式化验证的能力内核（金标准差距） & 机器检查证明 & klein2009 & \STRONG & ConsistenCy 无证明——正是 \S\ref{ch:limits} 的局限 1\\
撤销：永久单调 + 经中介 & 撤销理论 & redell1974; miller2003; miller2006 & \DIRECT & 退化（更易）情形；A4/A5 代码可证\\
vat/事件循环响应式并发 + 访问控制统一 & 健壮组合 & miller2006 & \DIRECT & 与“事件循环内核 + 响应式 Fiber”最同构\\
LLM 作为不可信提议者 & 混淆代理人/POLA & hardy1988; miller2003 & \DIRECT & 提示注入 $=$ 现代混淆代理人\\
能力式 Agent 控制层（平行现代工作） & 提示注入的系统性防御 & debenedetti2025 & \STRONG & 未经评审预印本；与 CaMeL 论点平行\\
intent 门（不可逆效果不内联） & 幂等意图/效果分级 & debenedetti2025 & \STRONG & 无已确立的“预算化权威”理论（诚实空白）\\
Fiber 监督生命周期 & Actor 语义/OTP 监督树 & agha1986; armstrong2003 & \STRONG & 分歧点：单次尝试诚实失败 vs OTP 重启掩盖\\
ACB 12 状态机 & 带标迁移系统 & keller1976 & \DIRECT & 字面实例\\
迁移不变式/fail-closed & 时序逻辑 & pnueli1977 & \DIRECT & S1--S3 的规范语言\\
内核为持久响应式运行时 & 响应式系统 & harel1985 & \STRONG & 响应式/变换式二分定位\\
积 LTS 的自动检验（未接线） & CTL 模型检验 & clarke1981 & \PROPOSED & 48 状态积空间是现成对象\\
工作流 DAG 校验 & WF-net 健全性 & vanderaalst1998 & \NA-for-now & 负结果：无环 DAG 上 Kahn 已给终止/无死锁\\
准入队列词汇/稳定性 & Kendall 记号/M/M/c & kendall1953; kleinrock1975 & \DIRECT & 词汇直用；公式是声明近似\\
积压-逗留不变式 & Little 定律 & little1961 & \DIRECT & 纪律无关——违背表下仍可用\\
近容量延迟爆炸 $\times$ 变异性 & Kingman 重交通 & kingman1962 & \DIRECT & GI/G/1 警示随行\\
优先级-then-FIFO 饥饿形状 & 非抢占 HOL & kleinrock1975 & \DIRECT & P1 截止期取消是兜底而非公平份额\\
公平份额替代（未实现） & 公平排队 & demers1989 & \STRONG & 对照框架\\
预算约束准入（研究模型） & CMDP & altman1999 & \PROPOSED & 现状是可行启发式；预算在授权处执法\\
LLM 服务时间非指数实证 & 迭代级调度/连续组批/内存管理 & yu2022orca; kwon2023vllm & \STRONG & $\mu$ 是客户端不可见的隐变量\\
发现-证据 why/where 溯源 & 数据溯源 & buneman2001 & \STRONG & 两个溯源问题逐字迁移\\
prov 四元组结构 & PROV-DM & provdm2013 & \DIRECT & 实体/活动/派生词汇的实例化\\
SHA-256 规范元组身份；快照寻址 & 内容寻址承诺/哈希树 & merkle1989 & \DIRECT & 单叶退化 Merkle\\
仅追加审计日志 & 安全审计日志/防篡改结构 & schneier1999; crosby2009 & \STRONG & 对手模型更弱；一致性证明是升级路径\\
分析器-步骤-材料-产品链 & 供应链证明 & torresarias2019 & \STRONG & in-toto 同构\\
模型输出事后接地 & 事后归因修订 & gao2023rarr & \STRONG & 接地依据是内核证据而非另一模型\\
原始模型输出不可靠的动机 & 审查评论处置率实证 & goldman2025 & \STRONG & ASE 2025 实证\\
神经代码审查基线（任务锚点） & 预训练代码审查 & li2022codereviewer & \STRONG & 无接地语义——正是差距\\
变更触发 vs 周期触发（结构） & 事件触发采样/最小事件间隔 & astrom1999; tabuada2007 & \STRONG(结构) & 定量最优性 \WEAK（无随机被控对象）\\
单运行去重 + 新到替换 & AoI 缓冲-1 洞察 & kaul2012 & \USEFUL & AoI 真正迁移的一条\\
AoI 词汇/峰值年龄记账 & AoI 综述 & yates2021 & \USEFUL & 词汇与记账框架\\
版本滞后 $D(t)$ & 版本 AoI & buyukates2021 & \USEFUL{} + \PROPOSED & 定义接地；Git-DAG/摘要等价类是扩展\\
摘要等值触发（无噪声观测） & 顺序变化检测 & page1954 & \NA & 负结果：无虚警/延迟权衡可优化\\
平均检测时延框架 & 更新报酬 & ross1970 & \USEFUL & 时延半边成立；代价半边未建模\\
Agent 框架对照（无权威中介） & ReAct / AutoGen / MetaGPT & yao2023react; wu2024autogen; hong2024metagpt & 对照 & 工作负载形态而非同类竞争者\\
\bottomrule
\caption{完整文献映射矩阵。强度列的分级记号（\DIRECT、\STRONG、\USEFUL、\PROPOSED、\WEAK、\NA）与正文六类分级一致。}
\label{tab:mapping}
\end{longtable}

% 附录 C：符号表
\chapter{符号表}\label{app:symbols}

\small
\begin{longtable}{@{}>{\raggedright\arraybackslash}p{3.4cm}>{\raggedright\arraybackslash}p{11.4cm}@{}}
\toprule
\textbf{符号} & \textbf{定义（一行）}\\
\midrule
\endhead
$\Nat,\;t$ & 离散时间域；步指标（内核中介循环的串行化序）\\
$X_t$ & 全局状态 $(R_t,S_t,Q_t,\mathfrak{A}_t,C_t,V_t,K_t,E_t,P_t,B_t)$\\
$R_t$ & 外部仓库状态（HEAD、对象库、工作树）——经观测摘要部分可观测\\
$S_t$ & 不可变钉定快照集 $(repo,sha)$\\
$Q_t$ & 准入队列（优先级 + FIFO 序）\\
$\mathfrak{A}_t$ & ACB 注册表：AgentId $\to$ 12 状态 LTS 状态\\
$C_t$ & 能力表：句柄 $\to$ 记录（subject/runId/action/resource/scope/expiresAt/revoked/budgetRef）\\
$V_t$ & 上下文可见性：内容寻址存储 $H:\text{text}\to\{0,1\}^{256}$ 上的逐 Agent 页表 $\mathrm{PT}_a$\\
$K_t$ & 共效/Fiber 可用性状态\\
$E_t$ & 证据库（不可变记录，内容派生指纹）\\
$P_t$ & 提供商状态（密钥、可用性、有效速率）——隐藏\\
$B_t$ & 逐能力预算态 $(used_c,reserved_c)$，不变式 $used+reserved\le max$\\
$\Pi,\;\Pi_{\mathrm{intent}}$ & 受保护系统调用动作集（12 项）；意图子集 $\{\tcode{repo.write},\tcode{github.publish}\}$\\
$q=(h,a,\alpha,r,s)$ & 系统调用请求：句柄、主语、动作、资源、作用域\\
$\kappa(r)$ & 资源类别函数（$\chi_5$ 使用）\\
$\chi_1..\chi_9$ & 九个授权检查（定义~\ref{def:auth}）\\
$\Auth$ & 授权谓词 $\bigwedge_i\chi_i$（有序首败拒绝）\\
$\Exec,\;m(q),\;\prec_{\mathrm{ser}}$ & 执行事件；其中介点；内核中介循环串行化序\\
S1, S2, S3 & 安全不变式（\S\ref{ch:security}、\S\ref{sec:revocation}、命题~\ref{prop:verdictplane}）\\
$t_r,\;t_{\mathrm{unload}},\;\delta_{\mathrm{prop}}$ & 撤销时刻；Fiber 卸载完成时刻；传播延迟\\
$W_{\mathrm{stale}}$ & 陈旧门面窗口 $[t_r,t_{\mathrm{unload}}]$\\
$\Avail,\;\Vis$ & 共效/Fiber 可用性；工作集页可见性\\
$\States,\;\to_A,\;T_{\mathrm{term}}$ & 12 状态 ACB LTS、其迁移关系、吸收终态集\\
$S_F,\;\to_F$ & 4 状态 Fiber LTS（Cordis 库语义）\\
$\varphi,\;\parallel_\varphi$ & 单向调度器$\to$Fiber 耦合；受限异步积\\
$G=(V,E),\;\sigma$ & 工作流 DAG；确定性拓扑序\\
$\mathcal{D}_t,\;Y_t,\;\varphi(\text{数据谓词})$ & 数据面坐标；权威坐标；数据面命题（重载 $\varphi$ 由语境区分）\\
$\mathcal{V}$ & 裁决泛函 $(q,X)\mapsto\Auth(q)$（命题~\ref{prop:verdictplane}）\\
$\operatorname{validate}$, $\operatorname{compile}$, $\operatorname{dryload}$, $\operatorname{execute}$ & 工作流四段管线函数\\
$e,\;\mathrm{prov},\;N,\;\canon,\;\fp$ & 证据记录；溯源；store 规范化；规范 JSON；SHA-256 指纹\\
$A_v$ & 版本钉定的全函数确定性分析器\\
$\ground$ & 确定性接地分类器 $\mathrm{Finding}^*\to\{\text{accept},\text{downgrade},\text{reject}\}^*$\\
$\lambda,\;\mu,\;\tau,\;c,\;\rho,\;A_e$ & 到达率、服务率、均值服务、并发上限、利用率、 offered load（erlangs）\\
$P_{\mathrm{wait}},\;W_q,\;W,\;L,\;L_q$ & Erlang-C 等待概率；排队/总逗留；Little 关系\\
$c_a^2,\;c_s^2$ & 到达间隔/服务时间的平方变异系数（Kingman）\\
$R_k,\;W_0$ & 类 $\le k$ 累积利用率；平均残余工作（HOL）——注意与仓库状态 $R_t$ 区分（下标）\\
$\mathcal{S},\mathcal{A},P,c,g_k,B_k,\gamma,\pi$ & CMDP 组元；折扣；策略\\
$\Delta(t),\;D(t),\;d(\cdot,\cdot),\;\Delta_v(t)$ & 时间年龄（已拒绝）；版本滞后；DAG 拟距离；版本 AoI\\
$T_{\mathrm{poll}},\;\delta_{\mathrm{scan}},\;T_{\mathrm{deb}},\;T_{\mathrm{run}}$ & 轮询周期、扫描时长、防抖（驻留时间）、运行时长；$W=\sum$（尾随窗口）\\
$\theta,\;\lambda_c$ & 变更时刻；泊松变更率（A10）\\
$\Phi,\;u_t,\;\zeta_t$ & 统一模型演化映射；控制；外生噪声（\PROPOSED）\\
A1--A10 & 常设假设（A1--A3 中介、A4--A5 撤销、A6--A8 分析器/规范、A7$'$ 单射、A9--A10 新鲜度）\\
F1--F5 & 运行活性的公平/终止假设（命题~\ref{prop:execsound}）\\
\bottomrule
\caption{符号表。重载符号（$R_k$ vs $R_t$、数据谓词 $\varphi$ vs 耦合 $\varphi$）已由上下文与下标消歧；本表为准。}
\label{tab:symbols}
\end{longtable}


\backmatter
\printbibliography[title={参考文献}]

\end{document}
